\section{Learning without pairs in blood: protocol, identification analyses and complete results}\label{sec:cross}
The protocol was written and registered
before extraction, fitting and prediction, and before any recipient scoring
cell was read. It is retrospective with respect to the released paper. The
colon cohort had informed method development; the blood recipient cohort
\cite{hao2021} had not been used in this project.

\subsection*{Data and features}
\textbf{Training populations.} Stephenson et al.\ \cite{stephenson2021}
(E-MTAB-10026): samples without
lipopolysaccharide stimulation; cells annotated as platelets or red cells
removed; patients as units (\XsPatients{} patients, 13 with two samples
pooled); at most 2,000 cells per patient, chosen pseudo-randomly
(\XsTrainCells{} cells). Within each patient a second pseudo-random split divided the cells
into an RNA half and a protein half (\XsHalfCells{} RNA-half cells); the
pairing-free estimators saw only RNA moments of the first and protein moments
of the second. Paired comparators used all selected cells.

\textbf{Recipients.} The confirmatory blood cohort was the 3$'$ CITE-seq study of
Hao et al.\ \cite{hao2021} (GEO GSE164378; \XsHaoDonors{} donors, three
vaccination time points pooled per donor, \XsHaoCells{} cells,
\XsHaoScoring{} scoring cells). The colon cohort was the 12 donors of the
large-panel analysis with a new pseudo-random split (\XsColonScoring{} scoring cells).
Scoring halves were written to separate files at split time and read only
after the predictions had been recorded.

\textbf{Features.} Proteins: the \XsProteins{} antibody targets present in all
three panels, matched by specificity (isotype controls and three ambiguous
targets excluded; where the blood recipient cohort measured a target with two
clones, the clone annotated as human-specific, otherwise the first, was used),
as $\log(1+\text{count})$ minus its mean over the matched targets of the
cell. Genes: $\log(1+{}$counts per $10^4)$; the 200 genes with the highest
variance-to-mean ratio over the selected training cells (mean above 0.05),
among symbols shared by all three studies and detected in at least 1\% of the
adaptation cells of both recipient studies, plus six panel genes not already
selected (\XsGenes{} genes). No protein value, pairing, scoring cell or
prediction entered feature selection.

\subsection*{Arms and tuning}
Pairing-free: (1) ridge regression of patient protein means on patient RNA
means in pooled within-patient standard-deviation units, penalty scale
\XsPfScale{} times the trace of the between-patient RNA-mean covariance,
interior to its grid; (2) the Gaussian moment likelihood of
Proposition~\ref{prop:pairing} with $W$ of rank \XsPfRank{}, penalty $10^{-3}$
times that scale, rank chosen from $\{2,4,8,16,32\}$ by held-out patients'
likelihood; the final fit stopped after \XsPfMomentIter{} iterations with
gradient norm $\XsPfMomentGrad$. Paired: the profile-likelihood interaction
(ridge \XsPairedRidge{}, interior), reference regression (penalty
\XsRegLambda{} times the mean diagonal, interior), pooled within-patient
cross-correlations and cross-covariances, and the closed form with identity
marginals. Benchmark: ridge regression on the recipient's own adaptation
pairs, tuned by three folds of adaptation cells. All tuned arms used a
five-value grid, three patient-grouped folds (patients assigned pseudo-randomly) and the
edge rule of the protocol.

\textbf{Amendment.} Before any prediction was computed, the rank search of
the moment estimator was restricted to its five-value grid without the edge
extension, because extending it to full rank (122) would have required three
cross-validation fits of more than an hour each on the available two cores;
the within-colon re-analysis kept the rank of its first version (8). The
information available at the time was the tuning outcome of the other arms and
the fact that the rank search had reached the edge of its grid. The chosen rank
(\XsPfRank) is therefore at the edge. The final fit stopped when the line
search could not decrease the objective further; three restarts from that point
changed the interaction by less than $\XsRestartBChange$ (relative) and no
recipient's endpoint by more than $\XsRestartEChange$.

\subsection*{Identifiability computations}
The effective supported subspace $S$ was spanned by the eigenvectors of the
between-patient RNA-mean matrix $G$ (pooled within-patient standard-deviation
units) whose ridge factor $g/(g+\lambda)$ was at least one half, with
$\lambda$ the fitted pairing-free penalty. $S$ is contained in the span $M$ of
Proposition~\ref{prop:means} and depends on the penalty. The share of the
recipient's standardized RNA variance inside $S$ is
$\operatorname{tr}(P_SR_xP_S)/\operatorname{tr}R_x$ for its adaptation
correlation matrix $R_x$ (the analyses of this note use measured $R_x$;
Supplementary Note~\ref{sec:noise} repeats them with its latent part). The
unsupported share is $\|R_x^{\mathrm{s}}(I-P_S)(R_x^{\mathrm{s}})^{-1}T\|_F/\|T\|_F$,
with $R_x^{\mathrm{s}}$ the shrunk RNA correlation matrix of the scoring cells
and $T$ their cross-correlation. The set of Proposition~\ref{prop:means}(4) was
computed with the adaptation correlation matrices and the fitted $W$ on $S$.
With measured RNA covariance, its whitened block $F$ had largest singular value
between \XsMaxFMin{} and \XsMaxFMax. It exceeded one in \XsMaxFAboveOne{} of
\XsIdUnits{} recipient analyses, including every total analysis (minimum
\XsMaxFTotalMin), so the exact set was empty there. Supplementary
Note~\ref{sec:noise} shows that RNA sampling noise explains most of this and
separates the remaining incompatibility from sampling uncertainty. The radii
below cap the singular values of $F$ at one, which is the same as clipping the
negative eigenvalues of $I-F^\top F$ at zero; they describe a repaired set, not
the identified one, and are kept for comparison with the earlier version. The
per-gene analysis (each gene's row of the cross-correlation matrix) was added
after scoring and is descriptive. In the per-gene analysis, the unsupported
share of each gene's relationships had Spearman correlation
\XsGeneSpearmanUnidHao{} with its prediction error in the confirmatory cohort and
\XsGeneSpearmanUnidColon{} in colon, whereas the share of the gene's own axis
inside $S$ had \XsGeneSpearmanAxisHao{} and \XsGeneSpearmanAxisColon. The
unsupported part of the target was \XsHaoIdFloor{} of its norm (range
\XsHaoIdFloorMin--\XsHaoIdFloorMax) in the confirmatory cohort's total analysis,
\XsHaoWIdFloor{} within its level-1 types and \XsColonIdFloor{} in colon; the
capped radius was \XsHaoIdRadius, \XsHaoWIdRadius{} and \XsColonIdRadius{} of
the norm of the truth, respectively.

\subsection*{Within-colon re-analysis of pairing-free learning}
This re-analysis replaces an earlier version of the within-colon analysis, with the same recipients, genes, proteins, units,
halves and endpoint. Each training biopsy's RNA moments came from one pseudo-random half
of its cells and its protein moments from the other. Folds for tuning were
donor-grouped. The ridge penalty was chosen from $\{10^{-4},\dots,1\}$ times the
trace of the between-biopsy RNA-mean covariance with the edge rule. The moment
estimator kept rank 8 and was run for up to 20,000 L-BFGS-B iterations
(gradient tolerance $10^{-7}$); within this limit, \XsRigConverged{} of 12 fits converged. The
other \XsRigContinued{} were continued post hoc from their final state for up to
20,000 further iterations, after
which \XsRigContinueConverged{} of them had converged; recipient errors changed by at
most \XsRigContinueMaxChange{} and interactions by at most \XsRigContinueMaxWChange{}
(relative Frobenius norm). Supplementary Table~\ref{tab:rigor} reports the errors
at the limit. The permutation control
reassigned protein moments between biopsies of different donors only, 200
times for the primary estimator and 20 times for the moment estimator, with
the same fitting rules. Supplementary Table~\ref{tab:rigor} gives every recipient.

\begin{table}[h]
\centering
\caption{\textbf{Within-colon pairing-free learning, corrected re-analysis.}
Relative cross-covariance error $E$ for each recipient, learned from the other
donors' biopsies without paired cells; $p$ is the donor-aware permutation
$p$-value (200 permutations for means, 20 for rank 8); $\dim M$ is the number
of gene directions the ridge estimator kept by at least one half. An asterisk
marks a rank-8 fit that stopped without meeting the convergence criterion.}
\label{tab:rigor}
\small
% Generated by extension/cross_study/make_assets.py. Do not edit.
\begin{tabular}{@{}lrrrrrrr@{}}
\toprule
Recipient & Biopsies & Penalty & Means $E$ & $p$ & Rank-8 $E$ & $p$ & $\dim M$\\
\midrule
HS1 & 19 & 0.01 & 0.612 & 0.005 & 0.619 & 0.048 & 11\\
HS2 & 19 & 0.01 & 0.662 & 0.005 & 0.701 & 0.048 & 11\\
HS3 & 19 & 0.01 & 0.470 & 0.005 & 0.632 & 0.048 & 12\\
HS4 & 19 & 0.01 & 0.422 & 0.005 & 0.940 & 0.143 & 11\\
HS5 & 19 & 0.01 & 0.399 & 0.005 & 0.547$^*$ & 0.048 & 11\\
HS6 & 19 & 0.01 & 0.588 & 0.005 & 0.624 & 0.048 & 11\\
HS7 & 19 & 0.01 & 0.637 & 0.005 & 0.696$^*$ & 0.048 & 11\\
HS8 & 20 & 0.01 & 0.555 & 0.005 & 0.478$^*$ & 0.048 & 11\\
HS9 & 19 & 0.01 & 0.468 & 0.005 & 0.763 & 0.095 & 11\\
HS10 & 20 & 0.01 & 0.827 & 0.005 & 0.806 & 0.048 & 12\\
HS11 & 19 & 0.01 & 0.579 & 0.005 & 0.617 & 0.048 & 11\\
HS12 & 20 & 0.01 & 0.574 & 0.005 & 0.734 & 0.048 & 12\\
\bottomrule
\end{tabular}

\end{table}

\subsection*{Complete results}
Supplementary Table~\ref{tab:cross} gives every arm's donor-mean error. The null of H1 kept
each training site's patients but paired their RNA and protein moments at
random within site; it therefore retained any site-level batch alignment,
which makes it conservative.

\begin{table}[h]
\centering
\caption{\textbf{Cross-study transfer: donor-mean relative error of the predicted
cross-correlation matrix.} Interactions were learned in the training study and
applied unchanged; the own paired half uses the recipient's adaptation pairs,
which the setting does not permit. Type columns centre every cell by its
cell-type mean in both assays. The last column is the mean correlation between
predicted and held-out entries in colon (all cells).}
\label{tab:cross}
\small
% Generated by extension/cross_study/make_assets.py. Do not edit.
\begin{tabular}{@{}lcccccc@{}}
\toprule
 & \multicolumn{3}{c}{Confirmatory blood cohort (8 donors)} & \multicolumn{3}{c}{Colon (12 donors)}\\
\cmidrule(lr){2-4}\cmidrule(l){5-7}
Arm & All cells & 8 types & 31 types & All cells & 14 types & Entry corr.\\
\midrule
No pairs: means (primary) & 0.292 & 1.082 & 1.484 & 0.988 & 1.274 & 0.40\\
No pairs: moments & 0.389 & 1.038 & 1.435 & 0.907 & 1.227 & 0.46\\
Paired closed form & 0.134 & 0.528 & 0.868 & 0.719 & 1.000 & 0.69\\
Reference regression & 0.522 & 0.680 & 0.842 & 0.859 & 0.983 & 0.52\\
Transferred correlation & 0.630 & 2.212 & 4.044 & 1.132 & 2.225 & 0.42\\
Transferred covariance & 0.951 & 8.960 & 18.953 & 4.651 & 15.642 & 0.23\\
Variances only & 0.976 & 0.901 & 0.915 & 0.971 & 0.983 & 0.28\\
Own paired half & 0.120 & 0.348 & 0.541 & 0.539 & 0.936 & 0.84\\
Independence & 1.000 & 1.000 & 1.000 & 1.000 & 1.000 & 0.00\\
\bottomrule
\end{tabular}

\end{table}

\subsection*{Post hoc analyses}
Two exploratory analyses were appended to the protocol after scoring; their
text and every version of the
protocol are kept with the code. Neither changes any confirmatory result.

\textbf{Between-type and within-type parts (P2).} For each recipient donor, the
scoring cells of the kept level-1 (blood) or coarse (colon) types were used to
compute the between-type part of the cross-covariance (type means weighted by
type proportions) and the within-type part (pooled within-type
cross-covariance); their sum is the total cross-covariance of those cells. Both
parts were divided by the scoring cells' total standard deviations. The prespecified
primary channel $W$ predicted each part from the donor's adaptation cells as
the corresponding RNA covariance part, in the same units, times $W^\top$, with
no saturation, so the two predicted parts also add to the direct-channel
prediction of the total. Supplementary Table~\ref{tab:mixture} reports the relative
Frobenius norm of each part of the truth, the relative error and the entry
correlation of each predicted part, and, as a descriptive addition computed
after P2 and not part of its plan, the share
$\operatorname{tr}(P_MR P_M)/\operatorname{tr}R$ of each RNA covariance part $R$
inside $M$. The direct channel's prediction of the within-type part had
\XsMixHaoNormWithin{} times the norm of the truth in the confirmatory cohort. The
direct channel without saturation is cruder than the closed form: for the total
over the kept types its error was \XsMixHaoErrDirect{} in the confirmatory cohort,
against \XsHaoErrPfMeans{} for the prespecified closed-form prediction of all cells.

\begin{table}[h]
\centering
\caption{\textbf{Between-type and within-type parts of the cross-covariance.}
Donor means. Norm share, Frobenius norm of each part of the scoring cells'
cross-covariance relative to the total (the parts are not orthogonal, so the
shares need not sum to one). Relative error and entry correlation, the prespecified
pairing-free channel's prediction of each part. Share inside $S$, fraction of
the corresponding RNA covariance of the adaptation cells inside the effective
supported subspace (descriptive addition).}
\label{tab:mixture}
\small
% Generated by extension/cross_study/make_assets.py. Do not edit.
\begin{tabular}{@{}lcccccccc@{}}
\toprule
 & \multicolumn{2}{c}{Norm share} & \multicolumn{2}{c}{Relative error} & \multicolumn{2}{c}{Entry correlation} & \multicolumn{2}{c}{Share inside $S$}\\
\cmidrule(lr){2-3}\cmidrule(lr){4-5}\cmidrule(lr){6-7}\cmidrule(l){8-9}
Cohort & between & within & between & within & between & within & between & within\\
\midrule
Confirmatory blood (8 donors) & 0.97 & 0.11 & 0.71 & 1.92 & 0.75 & 0.15 & 0.88 & 0.12\\
Colon (12 donors) & 0.81 & 0.43 & 1.14 & 1.38 & 0.35 & 0.07 & 0.57 & 0.09\\
\bottomrule
\end{tabular}

\end{table}

\textbf{Training on variation within cell types (P1).} Training units were the
training donors' cell types in the training study's initial clustering, with at least 50 cells in each
pseudo-random half (\XsPoneUnits{} units). RNA moments came from the RNA half and
protein moments from the protein half of each unit, and each unit's means were
centred on the cell-weighted mean of its type across donors, leaving only
variation among donors within types. The primary estimator was refitted by its
prespecified rules, with donor-grouped folds. The supported subspace had \XsPoneDim{}
directions, against \XsDimM{} for the prespecified fit, and the unsupported part of
the level-1 within-type truth fell to \XsPoneFloor{} of its norm (from
\XsHaoWIdFloor). The within-type error was \XsPoneErr{} at level 1 and
\XsPoneErrLTwo{} at level 2, against \XsHaoWErrPfMeans{} and
\XsHaoLTwoErrPfMeans{} for the prespecified fit, and the total error was
\XsPoneErrTotal. The training labels are the original authors' annotations of
the training study, which used both assays; the analysis therefore assumes
that each assay carries its own cell-type labels, as annotated atlases do.

\textbf{Gene--protein pairs.} Genes whose protein product was measured were
matched to their antibody targets by gene symbol (21 pairs among the prespecified
features). Their donor-mean held-out and predicted within-type correlations
(level-1 types, confirmatory cohort) are descriptive and were computed after
scoring; for FCGR3A and CD16 they were \XsPairFcgraWTruth{} held out,
\XsPairFcgraWPaired{} for the paired closed form, \XsPairFcgraWBench{} for the
recipient's own paired half and \XsPairFcgraWPfMeans{} for the primary
pairing-free channel; for IGHD and IgD, \XsPairIghdWTruth, \XsPairIghdWPaired,
\XsPairIghdWBench{} and \XsPairIghdWPfMeans; for KLRB1 and CD161,
\XsPairKlrbWTruth, \XsPairKlrbWPaired, \XsPairKlrbWBench{} and
\XsPairKlrbWPfMeans; and for IL7R and CD127, \XsPairIlrWTruth,
\XsPairIlrWPaired, \XsPairIlrWBench{} and \XsPairIlrWPfMeans. Supplementary Fig.~\ref{sfig:entries} shows held-out against
predicted entries of the cross-correlation matrix between and within cell types.

\begin{figure}[htbp]
\centering
\resizebox{\textwidth}{!}{% Generated by extension/cross_study/make_assets.py. Do not edit.
\definecolor{donorcol}{HTML}{5D6670}
\definecolor{haocol}{HTML}{217D91}
\definecolor{coloncol}{HTML}{A64442}
\begin{tikzpicture}[font=\sffamily]
\begin{axis}[name=a,width=.34\textwidth,height=.34\textwidth,xmin=-0.7,xmax=1.0,ymin=-0.7,ymax=1.0,
 xlabel={Held-out correlation},ylabel={Predicted correlation},title={\textbf{a}\enspace All cells, no pairs},
 axis x line*=bottom,axis y line*=left,tick align=outside,xmajorgrids=true,grid style={gray!18},axis line style={gray!55},tick label style={font=\sffamily\scriptsize},label style={font=\sffamily\small},clip=true,title style={font=\sffamily\small,at={(0,1)},anchor=base west,yshift=5pt}]
\draw[black!35] (axis cs:-0.7,-0.7) -- (axis cs:1.0,1.0);
\addplot[only marks,mark=*,mark size=.35pt,color=donorcol!45] coordinates {
 (-0.069,-0.040) (-0.061,-0.072) (-0.045,-0.035) (-0.176,-0.167) (-0.065,-0.023) (0.126,0.139) (0.120,0.106) (-0.062,-0.060) (-0.061,-0.043) (0.021,0.017) (-0.056,-0.037) (0.016,0.022)
 (0.154,0.095) (0.116,0.009) (0.007,0.005) (0.008,-0.008) (0.402,0.389) (0.598,0.487) (0.271,0.208) (0.505,0.439) (0.425,0.393) (0.380,0.399) (-0.484,-0.445) (-0.025,-0.020)
 (-0.092,-0.096) (-0.028,-0.045) (0.245,0.124) (-0.027,-0.026) (0.124,0.095) (0.130,0.019) (-0.029,-0.016) (-0.142,-0.097) (0.338,0.278) (-0.046,-0.040) (0.010,0.010) (-0.063,-0.024)
 (-0.161,-0.215) (0.042,0.051) (0.164,0.135) (-0.155,-0.111) (-0.179,-0.057) (0.296,0.156) (-0.414,-0.410) (0.627,0.505) (-0.218,-0.319) (0.137,0.047) (-0.228,-0.195) (-0.191,-0.176)
 (-0.193,-0.172) (-0.568,-0.495) (-0.381,-0.338) (0.184,0.131) (-0.264,-0.194) (-0.148,-0.141) (-0.012,0.193) (-0.083,-0.145) (0.297,0.246) (-0.051,-0.082) (-0.020,0.042) (0.090,-0.026)
 (-0.004,0.018) (-0.247,-0.200) (-0.234,-0.190) (-0.129,-0.097) (-0.659,-0.543) (-0.206,-0.275) (0.596,0.550) (0.520,0.340) (-0.236,-0.100) (-0.056,-0.054) (-0.013,-0.029) (-0.015,0.000)
 (-0.009,-0.023) (0.007,0.036) (0.112,0.111) (0.016,-0.011) (-0.025,0.010) (-0.037,-0.016) (-0.049,-0.038) (-0.020,0.031) (-0.040,-0.004) (-0.042,-0.036) (-0.051,-0.049) (0.065,0.060)
 (-0.093,-0.033) (-0.236,-0.233) (-0.085,-0.086) (0.724,0.511) (0.122,0.048) (0.366,0.289) (0.195,0.130) (-0.063,-0.060) (-0.179,-0.099) (-0.116,-0.089) (0.726,0.516) (-0.038,-0.020)
 (-0.172,-0.073) (-0.000,0.172) (-0.086,-0.054) (-0.089,-0.050) (-0.091,-0.126) (-0.122,-0.132) (0.203,0.236) (-0.272,-0.271) (0.400,0.391) (-0.109,-0.131) (0.035,0.062) (0.109,0.110)
 (0.402,0.312) (0.063,0.061) (-0.209,-0.167) (0.512,0.341) (-0.295,-0.252) (0.012,0.014) (0.120,0.098) (-0.020,0.235) (-0.104,-0.113) (0.432,0.263) (-0.051,-0.153) (-0.032,0.017)
 (0.027,-0.009) (-0.003,0.027) (-0.028,0.019) (-0.033,-0.010) (-0.023,-0.018) (-0.033,0.003) (-0.034,-0.052) (0.168,0.073) (-0.093,-0.031) (0.328,0.232) (-0.112,-0.122) (-0.023,-0.012)
 (-0.019,-0.031) (-0.025,-0.019) (-0.004,0.045) (0.141,0.181) (-0.000,-0.004) (-0.030,-0.057) (-0.128,-0.080) (-0.346,-0.315) (-0.221,-0.097) (-0.143,-0.095) (-0.145,-0.116) (-0.009,-0.021)
 (-0.109,-0.066) (-0.155,-0.116) (-0.213,-0.131) (0.010,0.069) (-0.113,-0.092) (-0.009,-0.086) (-0.008,-0.007) (-0.227,-0.145) (0.096,0.011) (-0.159,-0.167) (-0.014,-0.028) (-0.077,-0.085)
 (-0.015,-0.026) (0.018,0.020) (-0.014,0.064) (-0.027,-0.051) (-0.061,-0.079) (-0.093,-0.119) (-0.130,-0.146) (0.209,0.229) (-0.278,-0.284) (0.438,0.415) (-0.098,-0.145) (-0.009,0.032)
 (0.143,0.068) (-0.008,0.025) (0.039,0.036) (0.136,0.082) (0.214,0.130) (0.123,0.016) (-0.005,0.058) (0.062,0.043) (0.020,0.016) (-0.031,0.001) (-0.045,-0.069) (0.004,-0.096)
 (0.028,0.079) (-0.055,-0.150) (0.027,0.043) (-0.144,-0.134) (-0.143,-0.156) (-0.060,-0.083) (-0.408,-0.394) (-0.147,-0.134) (0.349,0.342) (0.281,0.321) (-0.134,-0.121) (-0.310,-0.209)
 (0.102,0.078) (-0.218,-0.128) (0.134,0.089) (0.317,0.279) (0.154,0.176) (0.063,0.058) (0.090,0.171) (0.368,0.379) (0.574,0.470) (0.225,0.221) (0.454,0.439) (0.400,0.378)
 (0.356,0.385) (-0.445,-0.430) (0.378,0.377) (0.769,0.637) (-0.099,-0.106) (-0.200,-0.123) (0.486,0.487) (-0.195,-0.136) (0.319,0.353) (-0.280,-0.180) (-0.491,-0.451) (-0.334,-0.277)
 (-0.268,-0.138) (-0.056,-0.039) (0.157,0.092) (-0.033,-0.150) (-0.195,-0.173) (-0.340,-0.283) (-0.069,-0.059) (0.246,0.201) (-0.291,-0.200) (-0.222,-0.139) (-0.207,-0.215) (-0.107,-0.077)
 (-0.132,-0.050) (-0.105,-0.112) (-0.022,-0.001) (-0.080,-0.093) (-0.276,-0.294) (-0.090,-0.114) (-0.024,0.043) (-0.108,-0.133) (-0.049,-0.062) (-0.035,0.027) (-0.118,-0.092) (0.594,0.429)
 (-0.082,-0.214) (-0.049,-0.037) (-0.008,0.015) (-0.021,0.010) (-0.001,-0.017) (-0.010,-0.044) (0.006,-0.003) (-0.005,-0.063) (-0.015,-0.078) (-0.040,-0.013) (0.063,0.165) (-0.085,-0.128)
 (0.557,0.421) (-0.232,-0.184) (0.402,0.326) (-0.244,-0.195) (-0.290,-0.210) (-0.264,-0.177) (-0.219,-0.176) (-0.206,-0.297) (-0.154,-0.090) (0.140,0.084) (0.002,0.073) (-0.140,-0.042)
 (-0.087,-0.058) (-0.126,-0.099) (0.059,0.041) (-0.133,-0.090) (-0.176,-0.103) (-0.025,0.052) (-0.095,-0.089) (-0.001,-0.055) (-0.006,-0.014) (-0.221,-0.129) (0.080,0.003) (-0.166,-0.142)
 (0.237,0.211) (-0.036,-0.061) (0.011,0.008) (-0.053,0.016) (-0.162,-0.203) (0.051,0.039) (0.143,0.124) (-0.054,-0.071) (-0.072,-0.091) (0.105,0.162) (-0.143,-0.158) (0.225,0.235)
 (-0.071,-0.057) (0.003,0.038) (-0.147,-0.123) (-0.083,-0.069) (-0.156,-0.105) (-0.426,-0.371) (-0.202,-0.206) (0.032,0.041) (-0.217,-0.165) (-0.103,-0.082) (-0.115,-0.134) (0.346,0.318)
 (-0.178,-0.042) (-0.319,-0.321) (-0.169,-0.128) (0.164,0.137) (-0.193,-0.162) (-0.090,-0.096) (-0.082,-0.101) (-0.045,-0.052) (-0.251,-0.261) (-0.091,-0.099) (0.205,0.209) (0.180,0.223)
 (-0.085,-0.083) (0.764,0.613) (-0.220,-0.171) (0.354,0.362) (-0.209,-0.183) (-0.298,-0.228) (-0.443,-0.364) (-0.219,-0.179) (-0.167,-0.321) (-0.174,-0.117) (0.156,0.068) (-0.015,0.078)
 (-0.143,-0.062) (-0.121,-0.089) (-0.125,-0.117) (0.115,0.072) (0.350,0.298) (0.681,0.575) (-0.098,-0.055) (-0.282,-0.131) (0.576,0.483) (-0.207,-0.117) (0.296,0.305) (-0.267,-0.280)
 (0.203,0.139) (-0.015,0.060) (-0.034,-0.076) (0.032,0.057) (-0.031,0.001) (-0.103,-0.106) (0.092,0.084) (0.107,0.133) (0.018,0.023) (0.014,0.004) (0.088,0.099) (0.084,0.004)
 (-0.102,-0.063) (0.211,0.366) (0.064,0.097) (-0.073,-0.053) (-0.054,-0.022) (-0.047,-0.038) (-0.086,-0.034) (-0.141,-0.073) (-0.120,-0.150) (0.340,0.268) (-0.056,-0.032) (-0.114,-0.134)
 (0.349,0.311) (-0.123,-0.031) (-0.312,-0.348) (-0.179,-0.149) (0.138,0.121) (-0.210,-0.183) (0.028,0.019) (0.026,0.052) (0.018,0.029) (0.084,0.098) (0.025,0.084) (-0.098,-0.097)
 (-0.039,-0.043) (-0.021,-0.014) (0.570,0.485) (-0.165,-0.132) (0.255,0.277) (-0.161,-0.147) (-0.222,-0.158) (-0.326,-0.264) (-0.169,-0.140) (-0.131,-0.215) (0.450,0.329) (0.620,0.486)
 (0.313,0.115) (0.501,0.348) (0.468,0.379) (0.451,0.369) (-0.506,-0.410) (-0.125,-0.100) (-0.175,-0.109) (0.021,0.055) (-0.089,-0.057) (0.002,-0.049) (-0.011,-0.006) (-0.176,-0.118)
 (0.080,-0.045) (-0.066,0.003) (0.220,0.175) (-0.024,-0.007) (0.003,-0.002) (-0.038,-0.030) (-0.091,-0.065) (0.012,0.009) (0.093,0.065) (-0.192,-0.157) (-0.238,-0.147) (0.403,0.267)
 (-0.544,-0.496) (0.825,0.678) (-0.282,-0.294) (0.184,0.073) (0.118,0.120) (0.371,0.294) (0.070,0.067) (-0.159,-0.128) (0.480,0.330) (-0.249,-0.229) (0.018,0.026) (0.119,0.100)
 (0.047,0.137) (-0.079,-0.065) (-0.113,-0.103) (0.141,0.019) (0.060,0.076) (0.037,-0.071) (0.056,0.034) (0.102,0.061) (0.086,0.062) (0.053,0.043) (0.266,0.202) (0.078,0.115)
 (-0.338,-0.258) (-0.118,-0.052) (-0.065,-0.051) (-0.126,-0.118) (-0.036,-0.028) (0.040,0.007) (-0.070,-0.037) (-0.042,-0.025) (0.114,0.137) (-0.026,-0.016) (-0.052,-0.022) (0.114,0.010)
 (0.266,0.158) (0.144,-0.079) (0.115,-0.024) (0.112,0.077) (0.136,0.046) (-0.161,-0.091) (-0.075,-0.025) (-0.205,-0.206) (-0.085,-0.089) (0.697,0.490) (0.116,0.060) (0.349,0.275)
 (0.222,0.138) (-0.073,-0.077) (-0.062,-0.030) (0.073,0.114) (-0.017,0.002) (0.001,0.030) (-0.027,-0.019) (-0.093,-0.190) (0.003,0.034) (0.066,0.071) (-0.142,-0.103) (0.087,0.109)
 (-0.070,-0.066) (-0.397,-0.299) (0.200,0.108) (-0.246,-0.170) (-0.005,-0.018) (0.106,0.067) (0.165,0.160) (0.030,0.029) (-0.084,-0.037) (0.306,0.155) (-0.118,-0.101) (0.008,0.040)
 (0.059,0.041) (-0.128,-0.131) (0.329,0.279) (-0.068,0.032) (-0.313,-0.335) (-0.185,-0.146) (0.229,0.169) (-0.202,-0.169) (-0.048,-0.027) (-0.046,-0.035) (-0.032,-0.046) (-0.036,-0.036)
 (-0.054,-0.079) (0.265,0.157) (-0.161,-0.024) (0.765,0.420) (-0.058,-0.008) (0.025,0.015) (-0.042,-0.040) (0.044,0.029) (0.062,0.059) (0.041,0.014) (0.011,0.001) (0.012,-0.014)
 (0.406,0.366) (0.564,0.442) (0.229,0.181) (0.467,0.415) (0.437,0.371) (0.401,0.385) (-0.456,-0.423) (-0.179,-0.140) (-0.233,-0.154) (0.063,0.090) (-0.111,-0.098) (-0.062,-0.120)
 (-0.004,-0.020) (-0.293,-0.199) (0.112,0.075) (-0.397,-0.369) (-0.275,-0.249) (-0.211,-0.140) (-0.038,-0.026) (0.128,0.104) (-0.088,-0.063) (-0.159,-0.142) (-0.275,-0.236) (-0.144,-0.126)
 (-0.173,-0.122) (0.273,0.218) (-0.402,-0.384) (0.638,0.559) (-0.211,-0.306) (0.050,0.039) (-0.074,-0.083) (-0.073,-0.029) (-0.076,-0.064) (-0.176,-0.157) (-0.126,-0.104) (0.058,0.099)
 (-0.089,-0.077) (-0.046,-0.045) (0.003,-0.137) (-0.018,0.009) (0.002,0.007) (0.021,-0.003) (0.011,-0.010) (-0.030,0.096) (0.013,0.002) (-0.165,-0.152) (-0.163,-0.166) (-0.076,-0.087)
 (-0.470,-0.433) (-0.170,-0.126) (0.402,0.386) (0.322,0.317) (-0.165,-0.137) (0.130,0.017) (-0.069,-0.040) (0.075,0.059) (-0.056,-0.030) (-0.080,-0.036) (-0.005,0.027) (-0.065,-0.025)
 (-0.075,-0.268) (-0.152,-0.090) (-0.366,-0.325) (-0.259,-0.099) (-0.176,-0.096) (-0.149,-0.125) (-0.017,-0.025) (-0.125,-0.072) (0.288,0.184) (0.494,0.392) (-0.091,-0.053) (-0.138,-0.023)
 (0.386,0.314) (-0.136,-0.053) (0.253,0.242) (-0.209,-0.151) (-0.377,-0.376) (-0.262,-0.236) (-0.213,-0.151) (-0.035,-0.034) (0.120,0.108) (-0.047,-0.010) (-0.148,-0.144) (-0.274,-0.238)
 (0.024,0.044) (0.035,0.026) (0.451,0.291) (-0.202,-0.094) (-0.215,-0.171) (0.292,0.244) (0.242,0.133) (0.074,0.081) (0.321,0.227) (0.040,0.040) (-0.112,-0.073) (0.368,0.202)
 (-0.242,-0.152) (0.022,0.038) (0.146,0.068) (-0.072,0.035) (0.263,0.190) (-0.084,-0.044) (-0.191,-0.247) (-0.107,-0.087) (0.057,0.086) (-0.128,-0.136) (0.036,0.029) (0.034,0.002)
 (0.025,-0.013) (0.108,0.057) (0.030,-0.028) (-0.116,-0.097) (-0.041,0.059) (-0.026,-0.031) (-0.056,-0.068) (0.018,0.036) (-0.015,-0.037) (0.006,0.017) (0.189,0.127) (0.062,0.066)
 (0.001,0.013) (0.015,0.052) (0.462,0.384) (0.601,0.468) (0.270,0.142) (0.523,0.404) (0.469,0.401) (0.468,0.411) (-0.494,-0.437) (0.205,0.051) (0.044,0.066) (-0.014,-0.019)
 (-0.023,-0.050) (-0.085,0.027) (0.018,-0.011) (0.123,-0.011) (-0.038,-0.050) (-0.469,-0.391) (-0.309,-0.240) (0.121,0.132) (-0.074,-0.017) (0.053,0.039) (-0.053,-0.047) (-0.185,-0.138)
 (-0.285,-0.245) (-0.193,-0.165) (0.018,0.031) (0.092,0.099) (-0.568,-0.485) (0.483,0.386) (-0.286,-0.262) (0.016,0.041) (-0.018,-0.003) (-0.020,0.022) (-0.013,-0.011) (-0.017,0.023)
 (-0.055,-0.016) (-0.021,-0.100) (0.063,0.063) (-0.016,-0.022) (-0.140,0.005) (0.207,0.132) (0.356,0.336) (-0.303,-0.404) (-0.197,-0.127) (0.316,0.172) (-0.179,-0.132) (-0.009,0.027)
 (-0.007,0.012) (-0.006,0.004) (0.012,0.016) (-0.004,-0.050) (0.064,0.032) (-0.062,-0.063) (0.174,0.120) (0.535,0.426) (-0.148,-0.125) (0.259,0.258) (-0.153,-0.135) (-0.184,-0.151)
 (-0.309,-0.242) (-0.154,-0.127) (-0.112,-0.227) (-0.166,-0.101) (-0.399,-0.352) (-0.285,-0.115) (-0.194,-0.109) (-0.159,-0.139) (-0.016,-0.026) (-0.126,-0.071) (0.321,0.206) (0.508,0.405)
 (-0.107,-0.050) (-0.139,-0.010) (0.456,0.366) (-0.127,-0.039) (0.346,0.268) (-0.237,-0.226) (-0.478,-0.433) (-0.326,-0.280) (-0.255,-0.148) (-0.047,-0.032) (0.165,0.116) (-0.076,-0.088)
 (-0.190,-0.172) (-0.329,-0.280) (-0.053,-0.053) (0.200,0.155) (-0.225,-0.155) (-0.165,-0.116) (-0.153,-0.158) (-0.090,-0.043) (-0.097,-0.039) (-0.042,-0.018) (-0.048,-0.006) (-0.044,-0.021)
 (-0.062,0.005) (-0.123,-0.023) (-0.064,-0.111) (0.243,0.184) (-0.051,-0.020) (-0.117,-0.133) (0.370,0.316) (-0.109,-0.004) (-0.306,-0.346) (-0.176,-0.156) (0.166,0.136) (-0.203,-0.180)
 (0.005,-0.001) (0.010,0.026) (-0.003,0.002) (0.039,0.054) (0.130,0.002) (-0.021,-0.022) (-0.066,-0.060) (0.080,0.115) (0.056,-0.029) (-0.042,-0.000) (0.068,0.040) (-0.035,0.013)
 (-0.051,0.000) (-0.009,-0.032) (-0.025,0.002) (-0.035,-0.098) (0.315,0.258) (0.403,0.300) (0.177,0.088) (0.343,0.277) (0.330,0.263) (0.312,0.275) (-0.331,-0.287) (-0.061,-0.053)
 (-0.136,-0.139) (-0.054,-0.061) (0.420,0.320) (0.054,0.036) (0.216,0.177) (0.100,0.080) (-0.039,-0.041) (-0.120,-0.077) (-0.074,-0.065) (0.469,0.351) (-0.028,-0.028) (-0.113,-0.046)
 (0.013,0.113) (-0.054,-0.050) (-0.066,-0.043) (0.072,0.085) (0.037,0.035) (0.366,0.259) (0.043,0.011) (-0.337,-0.258) (0.364,0.378) (0.196,0.131) (-0.027,-0.031) (-0.063,-0.029)
 (-0.039,-0.030) (-0.073,0.004) (-0.147,-0.059) (-0.089,-0.111) (0.281,0.225) (-0.041,-0.027) (-0.076,-0.047) (0.226,0.146) (-0.102,-0.078) (-0.198,-0.310) (-0.118,-0.109) (0.062,0.080)
 (-0.137,-0.138) (0.048,0.009) (0.042,0.027) (0.030,0.017) (0.140,0.086) (0.033,0.027) (-0.154,-0.122) (-0.037,0.014) (-0.028,-0.023) (-0.161,-0.160) (-0.025,0.014) (0.012,-0.027)
 (-0.052,-0.014) (-0.028,-0.001) (0.162,0.168) (0.004,0.015) (-0.046,-0.022) (0.431,0.319) (0.540,0.413) (0.265,0.063) (0.484,0.320) (0.419,0.348) (0.443,0.360) (-0.438,-0.383)
 (0.346,0.292) (0.602,0.509) (-0.095,-0.080) (-0.173,-0.078) (0.452,0.410) (-0.153,-0.100) (0.283,0.315) (-0.248,-0.193) (-0.468,-0.423) (-0.316,-0.272) (-0.256,-0.140) (-0.044,-0.032)
 (0.153,0.104) (-0.056,-0.102) (-0.184,-0.162) (-0.323,-0.266) (-0.172,-0.167) (-0.235,-0.195) (0.437,0.357) (-0.531,-0.503) (0.776,0.678) (-0.205,-0.267) (0.093,0.084) (0.138,0.103)
 (0.329,0.263) (0.058,0.057) (-0.163,-0.095) (0.381,0.283) (-0.224,-0.176) (0.030,0.043) (0.098,0.087) (0.027,0.033) (-0.038,0.002) (-0.054,-0.019) (0.039,-0.025) (0.042,0.052)
 (0.023,-0.045) (0.022,0.030) (-0.222,-0.173) (-0.211,-0.167) (-0.117,-0.081) (-0.589,-0.480) (-0.205,-0.248) (0.539,0.495) (0.418,0.282) (-0.205,-0.088) (0.668,0.575) (-0.195,-0.162)
 (0.342,0.327) (-0.179,-0.178) (-0.275,-0.198) (-0.466,-0.346) (-0.204,-0.176) (-0.148,-0.178) (-0.111,-0.063) (0.122,0.091) (0.015,0.065) (-0.096,-0.015) (-0.052,-0.031) (-0.110,-0.061)
 (-0.005,-0.002) (0.228,0.137) (0.403,0.322) (-0.070,-0.049) (-0.110,-0.011) (0.304,0.265) (-0.109,-0.043) (0.188,0.196) (-0.171,-0.089) (-0.119,-0.042) (-0.075,-0.051) (0.357,0.284)
 (-0.028,-0.003) (-0.057,-0.032) (0.036,0.097) (-0.044,-0.025) (-0.048,-0.033) (0.020,0.026) (0.027,0.002) (0.297,0.255) (-0.059,-0.074) (-0.167,-0.115) (0.298,0.299) (0.206,0.143)
 (-0.061,-0.032) (-0.006,0.002) (-0.050,-0.026) (-0.145,-0.112) (-0.062,-0.050) (0.003,0.035) (-0.057,-0.069) (-0.033,-0.026) (-0.045,-0.091) (0.192,0.226) (-0.075,-0.041) (-0.140,-0.223)
 (-0.079,-0.092) (0.054,0.061) (-0.087,-0.109) (-0.059,-0.013) (-0.060,-0.024) (-0.054,-0.031) (-0.061,-0.017) (-0.070,-0.135) (0.300,0.170) (-0.121,-0.084) (0.555,0.418) (-0.108,-0.111)
 (-0.034,-0.022) (0.022,-0.004) (-0.046,-0.027) (-0.024,-0.010) (0.133,0.138) (-0.006,-0.003) (-0.043,-0.027) (-0.139,-0.076) (-0.357,-0.315) (-0.253,-0.077) (-0.172,-0.073) (-0.142,-0.110)
 (-0.008,-0.015) (-0.130,-0.079) (0.362,0.309) (0.691,0.555) (-0.073,-0.068) (-0.200,-0.064) (0.491,0.450) (-0.192,-0.096) (0.363,0.345) (-0.277,-0.271) (-0.431,-0.382) (-0.308,-0.254)
 (-0.248,-0.125) (-0.057,-0.040) (0.153,0.093) (-0.031,-0.091) (-0.174,-0.155) (-0.305,-0.257) (-0.165,-0.164) (-0.216,-0.174) (0.330,0.271) (-0.475,-0.449) (0.744,0.655) (-0.215,-0.305)
 (0.023,0.049) (-0.142,-0.122) (-0.119,-0.102) (-0.132,-0.109) (-0.344,-0.309) (-0.240,-0.223) (0.065,0.092) (-0.170,-0.140) (-0.094,-0.085) (-0.111,-0.159) (0.311,0.255) (-0.145,-0.045)
 (-0.291,-0.370) (-0.166,-0.153) (0.159,0.154) (-0.191,-0.177) (0.184,0.167) (0.183,0.153) (0.089,0.074) (0.429,0.389) (0.162,0.237) (-0.532,-0.455) (-0.316,-0.273) (-0.098,-0.113)
 (0.356,0.324) (-0.102,-0.105) (0.205,0.168) (-0.113,-0.098) (-0.152,-0.107) (-0.244,-0.207) (-0.115,-0.103) (-0.074,-0.140) (-0.134,-0.099) (0.086,0.071) (-0.031,0.038) (-0.125,-0.046)
 (-0.074,-0.059) (-0.095,-0.090) (0.076,0.082) (-0.084,-0.059) (-0.114,-0.066) (0.009,0.030) (-0.041,-0.052) (-0.068,-0.084) (-0.001,-0.021) (-0.111,-0.063) (0.049,-0.010) (-0.111,-0.068)
 (-0.044,-0.064) (0.500,0.356) (-0.031,-0.022) (-0.125,-0.043) (-0.006,0.107) (-0.066,-0.035) (-0.073,-0.037) (0.012,0.033) (0.018,0.019) (0.285,0.199) (-0.124,-0.065) (-0.142,-0.132)
 (0.172,0.093) (0.185,0.101) (0.041,0.046) (0.064,0.060) (0.027,0.026) (0.023,-0.020) (0.160,0.127) (0.015,-0.079) (-0.002,-0.009) (0.048,0.025) (-0.009,-0.004) (0.035,0.047)
 (-0.031,-0.010) (-0.042,-0.076) (-0.001,-0.044) (-0.016,0.048) (-0.017,-0.023) (0.039,0.039) (0.035,0.023) (0.024,0.027) (0.113,0.076) (0.026,0.006) (-0.126,-0.099) (-0.039,0.013)
 (-0.019,-0.011) (0.746,0.615) (-0.235,-0.199) (0.412,0.378) (-0.251,-0.211) (-0.322,-0.243) (-0.442,-0.343) (-0.242,-0.198) (-0.188,-0.267) (0.206,0.207) (0.237,0.219) (0.064,0.126)
 (0.222,0.259) (0.215,0.205) (0.255,0.244) (-0.304,-0.287) (0.393,0.337) (0.743,0.611) (-0.097,-0.079) (-0.223,-0.098) (0.542,0.490) (-0.199,-0.109) (0.372,0.353) (-0.298,-0.268)
 (-0.189,-0.124) (-0.097,-0.083) (0.455,0.359) (-0.034,-0.025) (-0.083,-0.044) (-0.013,0.127) (-0.081,-0.061) (-0.091,-0.070) (-0.037,-0.034) (0.159,0.138) (-0.208,-0.126) (-0.141,-0.101)
 (-0.139,-0.144) (-0.080,-0.005) (-0.092,-0.024) (-0.022,-0.058) (-0.084,-0.063) (-0.067,-0.065) (-0.223,-0.138) (-0.160,-0.115) (0.062,0.062) (-0.081,-0.056) (-0.064,-0.053) (-0.052,0.101)
 (0.295,0.378) (-0.041,-0.108) (-0.108,-0.120) (-0.072,-0.068) (0.105,0.082) (-0.073,-0.061) (-0.022,-0.011) (-0.030,-0.031) (-0.006,-0.010) (-0.042,-0.053) (-0.026,-0.132) (0.008,0.032)
 (0.075,0.078) (-0.084,-0.038) (-0.104,-0.094) (-0.024,-0.018) (0.035,0.000) (-0.045,-0.023) (0.001,-0.008) (0.114,0.110) (-0.017,-0.010) (-0.039,-0.048) (-0.090,-0.059) (0.103,0.051)
 (0.003,0.033) (-0.065,-0.030) (-0.053,-0.039) (-0.070,-0.064) (0.021,0.013) (-0.084,-0.091) (-0.099,-0.074) (0.001,0.037) (-0.060,-0.037) (-0.011,-0.041) (-0.004,0.007) (-0.132,-0.097)
 (0.052,0.047) (-0.175,-0.123) (-0.093,-0.056) (-0.082,-0.096) (-0.028,-0.018) (0.042,0.051) (0.021,-0.038) (-0.063,-0.044) (-0.104,-0.081) (-0.074,-0.077) (-0.090,-0.080) (0.127,0.115)
 (-0.188,-0.158) (0.304,0.293) (-0.117,-0.112) (-0.008,0.019) (-0.208,-0.164) (-0.172,-0.141) (-0.175,-0.136) (-0.462,-0.403) (-0.327,-0.277) (0.171,0.123) (-0.234,-0.137) (-0.134,-0.111)
 (-0.104,-0.100) (0.332,0.341) (-0.132,-0.040) (-0.271,-0.296) (-0.163,-0.132) (0.155,0.118) (-0.181,-0.166) (-0.045,-0.030) (-0.048,-0.032) (-0.032,-0.037) (-0.041,-0.027) (-0.054,-0.080)
 (0.244,0.146) (-0.141,-0.045) (0.542,0.370) (0.512,0.379) (-0.147,-0.093) (0.200,0.240) (-0.129,-0.108) (-0.185,-0.122) (-0.257,-0.203) (-0.147,-0.098) (-0.118,-0.267) (-0.092,-0.065)
 (0.085,0.052) (-0.000,0.048) (-0.081,-0.018) (-0.056,-0.045) (-0.072,-0.063) (0.041,0.039) (-0.062,-0.008) (-0.073,-0.058) (-0.040,-0.009) (-0.007,-0.015) (0.044,0.009) (0.006,0.017)
 (-0.012,-0.014) (0.010,-0.052) (0.084,0.092) (0.059,0.061) (0.028,0.025) (0.024,0.033) (-0.000,-0.030) (-0.001,-0.018) (0.053,0.062) (0.071,0.065) (0.008,0.028) (0.014,0.000)
 (0.314,0.234) (-0.114,-0.085) (-0.111,-0.081) (0.208,0.190) (0.192,0.114) (-0.014,0.001) (0.091,0.088) (0.007,-0.003) (-0.066,-0.042) (0.121,0.098) (-0.166,-0.130) (0.143,0.131)
 (0.017,0.016) (-0.116,-0.146) (0.336,0.293) (-0.138,-0.017) (-0.297,-0.317) (-0.176,-0.143) (0.170,0.135) (-0.200,-0.175) (-0.160,-0.113) (-0.152,-0.122) (-0.079,-0.055) (-0.429,-0.344)
 (-0.153,-0.174) (0.394,0.335) (0.297,0.200) (-0.137,-0.108) (0.615,0.536) (-0.171,-0.132) (0.298,0.303) (-0.150,-0.153) (-0.147,-0.080) (-0.358,-0.283) (-0.193,-0.166) (-0.137,-0.138)
 (0.069,0.087) (0.113,0.083) (0.047,0.067) (0.096,0.128) (0.082,0.085) (0.061,0.061) (-0.073,-0.077) (-0.055,-0.027) (-0.148,-0.158) (-0.058,-0.065) (0.485,0.368) (0.104,0.055)
 (0.246,0.208) (0.151,0.099) (-0.049,-0.057) (-0.242,-0.208) (-0.157,-0.149) (0.322,0.277) (-0.033,-0.062) (-0.057,-0.008) (-0.018,0.157) (-0.099,-0.100) (-0.157,-0.126) (-0.093,-0.076)
 (0.103,0.129) (-0.116,-0.094) (-0.286,-0.206) (0.041,-0.022) (-0.153,-0.144) (-0.069,-0.030) (0.055,0.011) (0.094,0.076) (0.003,0.009) (-0.099,-0.057) (0.110,0.074) (-0.114,-0.092)
 (0.003,0.031) (0.039,0.016) (0.030,-0.060) (-0.062,-0.063) (-0.070,-0.060) (0.155,0.133) (0.044,0.028) (0.047,-0.009) (0.026,0.019) (-0.044,-0.015) (-0.047,-0.022) (-0.038,-0.031)
 (-0.028,-0.001) (-0.053,-0.081) (0.271,0.138) (-0.173,-0.088) (0.562,0.408) (-0.125,-0.104) (0.049,0.032) (-0.098,-0.066) (0.041,0.048) (0.370,0.205) (0.179,0.107) (0.015,0.015)
 (0.039,-0.080) (-0.070,-0.056) (0.026,0.006) (-0.018,0.035) (-0.068,-0.032) (-0.059,-0.041) (-0.053,-0.070) (0.069,0.048) (0.259,0.240) (0.480,0.413) (-0.062,-0.074) (-0.131,-0.077)
 (0.310,0.313) (-0.114,-0.088) (0.214,0.246) (-0.191,-0.122) (-0.124,-0.102) (-0.088,-0.055) (-0.071,-0.024) (-0.017,0.001) (0.062,-0.019) (-0.006,-0.050) (-0.050,-0.021) (-0.088,-0.072)
 (-0.041,-0.020) (-0.062,-0.091) (0.081,0.088) (-0.110,-0.071) (0.167,0.175) (-0.066,0.005) (0.053,0.015) (0.097,0.063) (0.066,0.074) (0.050,0.040) (0.023,0.019) (0.238,0.174)
 (0.028,-0.028) (0.002,0.025) (0.092,0.050) (-0.015,-0.014) (0.108,0.239) (-0.042,-0.067) (-0.052,-0.002) (-0.026,-0.046) (0.039,0.060) (-0.028,-0.043) (0.060,0.054) (0.047,0.046)
 (0.035,0.037) (0.186,0.159) (0.044,0.057) (-0.180,-0.146) (-0.079,-0.058) (0.026,0.051) (0.647,0.555) (-0.187,-0.166) (0.357,0.323) (-0.197,-0.180) (-0.266,-0.217) (-0.419,-0.333)
 (-0.203,-0.174) (-0.142,-0.196) (-0.058,-0.019) (0.019,-0.007) (-0.017,0.040) (-0.071,-0.015) (-0.032,-0.034) (-0.060,-0.041) (0.040,0.032) (-0.055,-0.058) (-0.103,-0.090) (-0.037,-0.005)
 (-0.030,-0.009) (-0.118,-0.127) (0.001,0.006) (-0.091,-0.049) (0.068,0.060) (-0.078,-0.041) (0.159,0.120) (-0.020,-0.018) (0.007,-0.011) (-0.030,0.002) (-0.044,-0.038) (0.012,0.009)
 (0.061,0.051) (-0.125,-0.100) (0.037,0.055) (0.015,0.011) (-0.384,-0.293) (0.234,0.162) (-0.171,-0.162) (-0.021,0.002) (0.027,0.039) (-0.028,-0.025) (0.061,0.056) (0.322,0.241)
 (-0.089,-0.044) (0.029,0.055) (0.027,0.019) (0.031,0.018) (0.039,-0.112) (-0.073,-0.063) (-0.079,-0.084) (0.066,0.115) (0.067,0.014) (0.033,0.020) (0.037,0.024) (-0.080,-0.049)
 (-0.072,-0.057) (-0.044,-0.034) (-0.203,-0.151) (-0.081,0.018) (0.195,0.172) (0.112,0.107) (0.006,0.006) (0.649,0.547) (-0.184,-0.172) (0.333,0.312) (-0.197,-0.172) (-0.260,-0.211)
 (-0.407,-0.335) (-0.193,-0.168) (-0.134,-0.212) (0.188,0.182) (0.229,0.225) (0.123,0.127) (0.205,0.206) (0.185,0.183) (0.186,0.199) (-0.183,-0.195) (0.363,0.344) (0.668,0.581)
 (-0.082,-0.094) (-0.185,-0.117) (0.472,0.461) (-0.165,-0.122) (0.305,0.320) (-0.275,-0.245) (-0.072,-0.063) (0.144,0.132) (-0.022,-0.012) (0.011,0.014) (0.008,-0.004) (-0.038,-0.139)
 (0.018,0.025) (0.083,0.068) (-0.177,-0.144) (-0.208,-0.135) (0.335,0.234) (-0.476,-0.442) (0.713,0.595) (-0.246,-0.250) (0.141,0.055) (0.108,0.030) (0.049,0.023) (0.037,0.024)
 (0.071,0.079) (-0.007,0.048) (0.067,0.047) (-0.006,-0.032) (0.016,0.014) (-0.097,-0.166) (0.293,0.272) (-0.136,-0.033) (-0.256,-0.319) (-0.146,-0.141) (0.116,0.122) (-0.175,-0.168)
 (-0.190,-0.156) (-0.173,-0.131) (-0.107,-0.084) (-0.429,-0.339) (-0.150,-0.249) (0.575,0.456) (0.173,0.137) (0.311,0.267) (-0.037,0.015) (0.017,0.040) (-0.028,-0.020) (0.029,0.056)
 (-0.018,-0.037) (-0.018,-0.091) (0.020,0.026) (0.013,0.006) (-0.048,-0.067) (-0.053,-0.070) (-0.023,-0.133) (-0.051,-0.123) (-0.048,-0.072) (-0.047,-0.076) (0.056,0.057) (0.038,0.073)
 (0.062,0.018) (-0.008,-0.021) (-0.030,-0.022) (0.039,0.008) (-0.031,0.006) (0.042,0.098) (-0.028,0.257) (-0.088,-0.060) (-0.050,-0.065) (0.377,0.289) (-0.015,-0.031) (-0.083,-0.036)
 (0.008,0.114) (-0.040,-0.045) (-0.048,-0.042) (-0.076,-0.030) (-0.083,-0.016) (0.137,0.033) (-0.156,-0.119) (0.247,0.170) (-0.138,-0.052) (0.167,0.039) (-0.033,-0.050) (-0.044,-0.025)
 (-0.034,-0.043) (-0.091,-0.114) (-0.079,-0.066) (0.057,0.033) (-0.051,-0.055) (-0.021,-0.031) (-0.062,0.020) (0.117,0.075) (0.154,0.167) (-0.164,-0.245) (-0.098,-0.071) (0.111,0.073)
 (-0.097,-0.071) (-0.131,-0.115) (-0.111,-0.106) (-0.072,-0.066) (-0.354,-0.295) (-0.084,-0.220) (0.285,0.292) (0.235,0.183) (-0.188,-0.090) (-0.362,-0.303) (0.153,0.112) (-0.250,-0.228)
 (0.124,0.127) (0.455,0.288) (0.389,0.224) (0.126,0.100) (0.127,0.026) (0.317,0.259) (0.407,0.302) (0.175,0.108) (0.347,0.280) (0.330,0.262) (0.299,0.283) (-0.332,-0.289)
 (0.206,0.128) (0.419,0.325) (-0.023,-0.033) (-0.165,-0.064) (0.228,0.261) (-0.139,-0.069) (0.154,0.208) (-0.149,-0.040) (-0.023,0.017) (0.104,0.094) (0.026,0.010) (0.002,-0.030)
 (-0.031,-0.006) (0.008,0.033) (0.017,0.005) (0.027,0.048) (0.031,0.077) (0.094,0.177) (-0.268,-0.251) (0.197,0.140) (-0.196,-0.260) (0.048,0.161) (-0.044,-0.069)};
\draw[haocol!60,line width=.3pt] (axis cs:0.890,0.766) -- (axis cs:0.350,-0.080);
\addplot[only marks,mark=*,mark size=1.6pt,color=haocol] coordinates { (0.890,0.766) };
\node[font=\sffamily\tiny,anchor=west,color=haocol,inner sep=1pt] at (axis cs:0.360,-0.080) {\textit{LYZ}--CD371};
\draw[haocol!60,line width=.3pt] (axis cs:0.869,0.609) -- (axis cs:0.350,-0.165);
\addplot[only marks,mark=*,mark size=1.6pt,color=haocol] coordinates { (0.869,0.609) };
\node[font=\sffamily\tiny,anchor=west,color=haocol,inner sep=1pt] at (axis cs:0.360,-0.165) {\textit{CD79A}--CD19};
\draw[coloncol!60,line width=.3pt] (axis cs:0.819,0.570) -- (axis cs:0.350,-0.250);
\addplot[only marks,mark=*,mark size=1.6pt,color=coloncol] coordinates { (0.819,0.570) };
\node[font=\sffamily\tiny,anchor=west,color=coloncol,inner sep=1pt] at (axis cs:0.360,-0.250) {\textit{IL7R}--CD127};
\draw[coloncol!60,line width=.3pt] (axis cs:0.765,0.420) -- (axis cs:0.350,-0.335);
\addplot[only marks,mark=*,mark size=1.6pt,color=coloncol] coordinates { (0.765,0.420) };
\node[font=\sffamily\tiny,anchor=west,color=coloncol,inner sep=1pt] at (axis cs:0.360,-0.335) {\textit{IGHD}--IgD};
\draw[coloncol!60,line width=.3pt] (axis cs:0.755,0.294) -- (axis cs:0.350,-0.420);
\addplot[only marks,mark=*,mark size=1.6pt,color=coloncol] coordinates { (0.755,0.294) };
\node[font=\sffamily\tiny,anchor=west,color=coloncol,inner sep=1pt] at (axis cs:0.360,-0.420) {\textit{FCGR3A}--CD16};
\draw[coloncol!60,line width=.3pt] (axis cs:0.688,0.259) -- (axis cs:0.350,-0.505);
\addplot[only marks,mark=*,mark size=1.6pt,color=coloncol] coordinates { (0.688,0.259) };
\node[font=\sffamily\tiny,anchor=west,color=coloncol,inner sep=1pt] at (axis cs:0.360,-0.505) {\textit{KLRB1}--CD161};
\end{axis}
\begin{axis}[at={(a.east)},anchor=west,xshift=1.3cm,name=b,width=.34\textwidth,height=.34\textwidth,xmin=-0.6,xmax=0.8,ymin=-0.6,ymax=0.8,
 xlabel={Held-out correlation},ylabel={Predicted correlation},title={\textbf{b}\enspace Within types, no pairs},
 axis x line*=bottom,axis y line*=left,tick align=outside,xmajorgrids=true,grid style={gray!18},axis line style={gray!55},tick label style={font=\sffamily\scriptsize},label style={font=\sffamily\small},clip=true,title style={font=\sffamily\small,at={(0,1)},anchor=base west,yshift=5pt}]
\draw[black!35] (axis cs:-0.6,-0.6) -- (axis cs:0.8,0.8);
\addplot[only marks,mark=*,mark size=.35pt,color=donorcol!45] coordinates {
 (-0.019,0.013) (-0.014,-0.027) (-0.022,-0.012) (-0.054,-0.042) (-0.025,0.021) (0.007,0.016) (0.005,0.004) (-0.008,-0.011) (-0.032,-0.037) (0.001,0.008) (-0.031,-0.020) (-0.001,0.005)
 (-0.025,-0.003) (0.018,-0.042) (0.001,0.004) (-0.006,-0.005) (-0.122,0.004) (0.080,0.030) (-0.024,0.106) (-0.040,0.041) (-0.030,0.009) (-0.155,-0.024) (-0.015,-0.025) (0.014,-0.013)
 (-0.008,-0.005) (0.007,-0.006) (-0.034,-0.072) (-0.095,-0.070) (-0.010,-0.013) (0.071,-0.016) (-0.008,0.017) (-0.224,-0.150) (-0.019,0.020) (-0.008,0.003) (-0.001,0.034) (-0.001,0.002)
 (-0.098,-0.291) (0.002,0.041) (-0.001,0.023) (-0.014,0.031) (-0.010,0.068) (0.025,-0.096) (-0.005,-0.043) (0.026,0.107) (-0.035,-0.159) (0.065,0.002) (-0.029,-0.032) (-0.005,-0.038)
 (0.008,-0.021) (-0.047,-0.015) (-0.023,-0.002) (0.014,-0.002) (0.002,-0.012) (-0.005,-0.023) (0.016,0.152) (-0.019,-0.066) (-0.186,-0.022) (0.002,0.024) (0.014,0.054) (0.052,-0.034)
 (0.006,0.008) (-0.025,0.018) (-0.017,0.005) (-0.021,-0.000) (-0.065,-0.015) (-0.057,-0.108) (0.083,0.037) (0.038,-0.042) (-0.009,0.007) (0.029,0.009) (-0.001,-0.020) (-0.023,0.005)
 (0.015,-0.011) (0.009,0.004) (0.039,0.036) (0.017,-0.004) (-0.006,0.028) (-0.002,0.006) (-0.000,-0.003) (-0.004,0.041) (0.002,0.027) (-0.004,-0.004) (-0.005,-0.017) (-0.001,0.006)
 (0.004,-0.002) (-0.016,-0.014) (0.005,-0.006) (0.078,0.076) (0.044,0.043) (0.021,0.039) (-0.065,-0.017) (-0.003,0.029) (0.006,-0.001) (-0.021,-0.020) (0.080,0.079) (0.001,-0.005)
 (-0.002,0.001) (0.023,0.064) (-0.005,-0.009) (-0.004,0.002) (0.004,-0.049) (-0.008,-0.030) (-0.008,0.089) (-0.006,-0.011) (0.048,0.077) (0.020,-0.013) (-0.036,-0.002) (-0.024,0.014)
 (0.078,0.068) (-0.017,0.020) (-0.459,-0.268) (-0.036,-0.011) (-0.010,-0.061) (-0.009,-0.000) (0.014,0.029) (0.005,0.178) (-0.005,-0.008) (-0.009,0.011) (-0.002,-0.113) (0.001,0.028)
 (0.001,-0.019) (-0.003,0.017) (-0.002,0.043) (-0.009,0.003) (-0.003,-0.001) (-0.023,0.024) (-0.002,-0.017) (0.007,-0.010) (0.003,0.007) (0.002,0.028) (-0.062,-0.081) (-0.002,0.001)
 (-0.060,-0.046) (0.012,-0.006) (0.016,0.032) (0.059,0.117) (0.009,0.000) (0.005,-0.006) (0.008,0.016) (-0.026,-0.059) (0.001,-0.004) (0.003,0.003) (-0.025,-0.014) (0.001,0.022)
 (-0.005,-0.015) (0.006,-0.004) (-0.030,0.005) (0.003,0.019) (-0.038,-0.035) (0.086,-0.052) (-0.004,0.015) (-0.008,-0.006) (-0.002,-0.017) (-0.049,-0.084) (0.011,-0.019) (0.000,0.005)
 (0.002,-0.003) (-0.002,0.004) (-0.014,0.061) (0.017,-0.011) (-0.009,-0.035) (0.009,-0.031) (-0.005,-0.037) (-0.001,0.059) (-0.001,-0.017) (0.058,0.031) (0.043,-0.002) (-0.093,-0.044)
 (0.062,0.023) (-0.069,-0.026) (-0.014,-0.011) (-0.030,-0.043) (-0.006,0.003) (0.171,0.076) (-0.016,0.037) (-0.023,-0.000) (0.001,-0.001) (-0.003,0.027) (0.004,-0.020) (-0.041,-0.164)
 (0.004,0.044) (-0.066,-0.123) (0.008,0.021) (-0.000,-0.003) (-0.008,-0.048) (0.017,-0.029) (-0.060,-0.108) (-0.022,0.001) (0.017,0.046) (-0.006,0.133) (-0.003,-0.038) (-0.037,-0.025)
 (-0.013,-0.002) (-0.016,0.039) (0.038,0.016) (-0.059,-0.022) (-0.234,-0.088) (-0.022,-0.003) (-0.020,0.069) (-0.100,0.017) (0.071,0.044) (-0.048,0.142) (-0.053,0.076) (-0.013,0.010)
 (-0.126,-0.015) (0.001,-0.034) (0.017,0.105) (0.239,0.107) (0.006,-0.044) (0.004,-0.004) (-0.075,0.049) (-0.033,-0.034) (-0.091,0.002) (-0.007,0.160) (-0.018,-0.042) (-0.014,0.006)
 (-0.030,-0.016) (-0.000,-0.001) (-0.023,-0.024) (0.034,-0.119) (-0.002,-0.006) (-0.052,-0.010) (-0.013,-0.018) (0.026,0.049) (-0.019,-0.007) (-0.034,-0.074) (-0.056,-0.114) (0.004,-0.058)
 (-0.008,0.027) (-0.018,-0.052) (-0.009,0.023) (-0.001,-0.022) (-0.030,-0.059) (-0.004,-0.007) (-0.063,0.020) (-0.002,-0.029) (-0.001,-0.019) (0.003,-0.074) (-0.006,0.031) (0.099,0.084)
 (-0.009,-0.150) (0.002,-0.025) (-0.083,0.006) (-0.017,0.006) (0.001,0.006) (0.002,-0.030) (-0.000,-0.007) (-0.037,-0.077) (-0.006,-0.083) (0.010,0.014) (0.018,0.119) (0.003,-0.028)
 (0.100,0.080) (-0.017,-0.007) (0.050,0.057) (0.019,0.009) (-0.014,-0.016) (-0.062,-0.054) (-0.013,-0.017) (-0.012,-0.165) (-0.003,0.008) (0.235,0.060) (0.092,0.078) (0.105,0.070)
 (0.017,0.001) (0.000,-0.005) (-0.110,-0.068) (0.005,0.020) (-0.006,0.018) (-0.020,0.005) (-0.019,-0.033) (0.049,-0.008) (-0.007,0.006) (-0.021,-0.007) (-0.005,-0.013) (-0.314,-0.232)
 (-0.004,-0.006) (0.001,-0.005) (-0.005,0.022) (0.002,0.038) (-0.092,-0.238) (0.002,0.023) (-0.000,0.026) (-0.002,-0.022) (-0.005,-0.016) (-0.005,0.059) (-0.002,-0.015) (0.025,0.035)
 (-0.001,0.012) (-0.031,-0.006) (-0.018,-0.012) (-0.023,0.004) (-0.031,-0.001) (0.004,-0.047) (-0.007,-0.041) (-0.025,-0.045) (-0.035,-0.049) (-0.014,-0.004) (-0.007,-0.038) (0.038,0.065)
 (0.008,-0.016) (-0.077,-0.031) (-0.005,0.018) (-0.010,-0.005) (-0.002,0.006) (0.002,-0.009) (0.006,-0.028) (0.005,-0.015) (-0.016,-0.063) (-0.013,-0.008) (-0.020,0.010) (-0.011,0.084)
 (0.001,-0.008) (0.124,0.100) (-0.022,0.010) (-0.036,0.047) (0.009,0.015) (-0.010,-0.005) (0.029,-0.055) (0.005,0.003) (-0.013,-0.251) (-0.006,-0.019) (0.352,0.094) (0.079,0.105)
 (0.150,0.078) (0.002,-0.023) (0.044,-0.003) (-0.096,-0.045) (0.022,0.026) (0.017,0.060) (0.007,0.019) (-0.013,-0.012) (0.077,0.016) (-0.010,0.001) (-0.038,0.010) (-0.011,-0.038)
 (0.035,-0.000) (-0.002,-0.015) (-0.005,-0.020) (0.001,0.010) (-0.001,-0.001) (-0.001,-0.007) (0.010,0.002) (-0.004,0.013) (-0.009,0.009) (-0.001,-0.006) (0.109,0.063) (0.014,-0.006)
 (0.006,0.023) (0.149,0.303) (0.051,0.058) (-0.008,-0.008) (-0.008,-0.011) (-0.001,-0.006) (-0.076,-0.055) (-0.007,-0.004) (-0.009,-0.026) (0.001,0.020) (-0.001,0.004) (0.022,-0.028)
 (0.048,0.086) (0.013,-0.036) (-0.003,-0.040) (0.006,-0.005) (-0.126,-0.105) (-0.011,-0.016) (-0.002,-0.009) (-0.000,0.027) (-0.003,0.009) (-0.005,0.039) (0.000,0.072) (-0.003,-0.024)
 (-0.008,-0.019) (-0.004,-0.015) (0.054,0.075) (-0.007,0.007) (-0.026,0.025) (-0.001,0.001) (-0.005,0.002) (0.028,0.006) (-0.004,0.004) (-0.013,-0.086) (0.048,0.019) (0.001,0.055)
 (0.023,-0.100) (0.040,-0.026) (0.046,0.059) (0.078,0.033) (-0.031,-0.005) (0.001,-0.014) (-0.030,-0.006) (0.016,0.013) (-0.009,-0.000) (0.056,0.010) (-0.009,0.007) (-0.001,0.008)
 (0.005,-0.095) (-0.011,0.049) (0.000,0.037) (0.002,0.002) (0.001,0.001) (0.001,-0.015) (-0.033,-0.033) (-0.000,-0.013) (0.006,0.003) (-0.002,0.029) (-0.010,0.016) (0.071,0.001)
 (-0.044,-0.045) (0.023,0.100) (-0.047,-0.017) (0.143,0.026) (-0.025,0.021) (0.045,0.050) (-0.019,0.018) (-0.421,-0.254) (-0.025,-0.016) (0.032,-0.037) (-0.010,0.011) (0.004,0.025)
 (-0.004,0.044) (0.003,0.009) (-0.005,0.007) (0.029,-0.039) (-0.003,0.002) (0.058,-0.013) (-0.003,-0.014) (-0.011,-0.028) (-0.015,-0.009) (-0.019,-0.009) (-0.057,-0.018) (-0.018,0.035)
 (-0.020,-0.021) (0.009,0.057) (-0.014,0.001) (-0.027,-0.015) (0.004,0.000) (0.007,0.004) (-0.007,0.002) (-0.023,-0.009) (-0.070,-0.021) (-0.018,-0.008) (0.002,-0.000) (0.111,0.007)
 (0.088,0.019) (0.100,-0.163) (0.102,-0.071) (0.045,0.042) (0.126,0.016) (-0.083,-0.028) (-0.002,-0.018) (-0.011,-0.007) (-0.000,-0.012) (0.075,0.055) (-0.010,0.025) (0.022,0.024)
 (-0.033,-0.018) (-0.007,0.017) (-0.060,-0.005) (-0.072,0.010) (-0.005,-0.004) (-0.004,0.031) (-0.004,-0.005) (-0.045,-0.163) (-0.012,0.023) (0.000,0.021) (-0.025,-0.004) (-0.008,0.023)
 (0.028,0.021) (-0.015,-0.062) (0.022,0.069) (-0.084,0.020) (0.095,0.030) (0.071,0.025) (-0.078,0.013) (-0.002,0.002) (-0.034,-0.008) (0.090,-0.001) (0.091,0.018) (-0.003,0.032)
 (0.011,-0.001) (-0.009,-0.040) (0.044,0.031) (0.004,-0.017) (-0.014,-0.021) (-0.012,-0.008) (0.099,0.055) (-0.007,-0.006) (-0.000,0.004) (0.001,-0.008) (0.005,-0.021) (0.010,-0.016)
 (0.003,-0.018) (-0.042,-0.000) (-0.015,0.091) (0.514,0.101) (0.006,0.022) (-0.000,-0.001) (-0.001,-0.008) (0.021,0.009) (-0.017,-0.021) (-0.034,-0.037) (-0.010,-0.011) (-0.009,-0.042)
 (-0.059,0.005) (0.042,0.026) (-0.043,0.074) (-0.037,0.049) (0.027,0.015) (-0.067,0.001) (-0.005,-0.032) (-0.002,-0.008) (-0.001,0.015) (0.047,0.053) (-0.027,-0.020) (-0.002,-0.081)
 (-0.001,-0.011) (-0.032,-0.059) (-0.005,0.016) (-0.000,-0.008) (-0.008,-0.032) (0.003,-0.025) (0.016,-0.003) (-0.005,0.005) (-0.072,0.010) (0.001,-0.009) (-0.022,-0.004) (-0.006,-0.007)
 (-0.008,-0.002) (-0.049,0.020) (-0.007,-0.014) (0.123,0.214) (-0.039,-0.162) (-0.070,-0.032) (-0.003,-0.024) (-0.003,0.028) (-0.012,-0.012) (-0.025,-0.041) (-0.007,-0.008) (-0.016,0.052)
 (-0.003,-0.001) (-0.005,-0.005) (-0.004,-0.155) (-0.003,0.010) (0.007,-0.017) (0.012,0.024) (-0.001,-0.024) (-0.007,0.130) (0.003,-0.004) (0.001,0.002) (-0.005,-0.043) (0.011,-0.026)
 (-0.054,-0.100) (-0.020,0.044) (-0.001,0.037) (-0.007,0.091) (-0.005,-0.028) (-0.042,-0.071) (-0.011,0.009) (-0.079,-0.021) (0.014,0.024) (0.012,0.024) (0.104,0.049) (-0.000,0.020)
 (-0.014,-0.207) (-0.006,0.004) (-0.024,-0.031) (-0.018,0.011) (-0.040,-0.003) (-0.009,-0.013) (-0.016,-0.006) (-0.010,-0.005) (0.020,-0.033) (0.026,0.070) (-0.021,0.002) (-0.002,0.013)
 (0.050,-0.008) (-0.013,0.013) (-0.000,0.014) (-0.015,-0.009) (-0.005,-0.046) (-0.009,-0.031) (-0.007,-0.028) (0.016,-0.009) (-0.012,0.015) (-0.031,0.069) (0.005,-0.015) (-0.042,-0.025)
 (-0.004,0.019) (0.006,0.011) (0.119,0.068) (-0.006,-0.057) (-0.000,-0.031) (0.047,0.034) (0.027,-0.003) (-0.008,0.021) (0.007,0.026) (0.006,0.009) (-0.005,-0.025) (0.112,-0.005)
 (0.031,0.003) (0.007,0.030) (0.063,0.013) (0.001,0.146) (0.088,0.008) (-0.003,-0.003) (-0.009,-0.070) (-0.000,-0.001) (-0.068,-0.019) (-0.008,-0.037) (-0.003,0.016) (-0.000,-0.016)
 (-0.003,-0.028) (-0.006,-0.017) (-0.004,-0.056) (0.006,-0.032) (0.011,0.085) (-0.003,-0.013) (0.002,-0.022) (-0.001,0.024) (0.015,-0.004) (-0.008,0.001) (0.062,0.012) (-0.045,0.003)
 (-0.007,0.007) (0.004,0.050) (0.022,0.043) (0.093,0.071) (0.006,-0.013) (0.036,0.028) (0.062,0.066) (0.034,0.036) (-0.020,-0.025) (0.067,0.004) (-0.017,0.017) (0.010,-0.016)
 (0.021,0.001) (-0.079,-0.016) (0.030,0.026) (-0.000,-0.041) (-0.015,0.019) (-0.058,-0.061) (-0.027,-0.014) (0.005,0.012) (-0.017,0.029) (-0.011,-0.003) (0.009,-0.048) (-0.005,0.022)
 (-0.030,-0.039) (-0.017,-0.015) (-0.002,0.011) (0.027,0.073) (-0.020,-0.074) (0.074,0.160) (-0.045,-0.007) (0.019,0.040) (-0.016,-0.003) (0.003,0.035) (0.006,0.002) (0.050,0.059)
 (0.008,-0.009) (0.002,-0.055) (-0.032,-0.016) (0.002,-0.004) (-0.022,0.052) (0.060,0.001) (0.042,0.074) (0.029,-0.093) (-0.024,0.001) (0.202,0.053) (-0.011,-0.002) (-0.005,0.033)
 (-0.003,0.015) (-0.001,0.009) (-0.001,0.011) (0.004,-0.032) (0.008,0.006) (-0.001,-0.028) (0.028,0.011) (0.070,0.010) (-0.005,-0.004) (0.006,0.022) (-0.006,-0.004) (0.000,-0.001)
 (-0.009,0.003) (-0.002,-0.001) (-0.004,-0.105) (-0.010,0.005) (-0.047,-0.041) (-0.031,-0.004) (-0.050,-0.016) (-0.008,-0.022) (-0.006,0.010) (-0.005,-0.002) (0.017,-0.046) (-0.124,-0.025)
 (-0.027,0.028) (-0.003,0.000) (0.075,0.002) (0.008,0.028) (0.074,0.001) (-0.013,-0.084) (-0.009,-0.008) (-0.015,-0.023) (-0.011,-0.025) (0.017,-0.000) (0.004,0.008) (-0.050,0.027)
 (-0.001,-0.023) (-0.041,-0.019) (-0.009,-0.016) (0.021,0.026) (-0.007,0.002) (-0.010,-0.060) (-0.011,-0.058) (-0.009,-0.005) (-0.000,0.015) (-0.010,0.011) (-0.001,0.017) (-0.003,0.007)
 (0.021,0.032) (-0.001,0.020) (-0.006,-0.013) (0.035,0.013) (-0.011,0.014) (0.006,-0.035) (0.122,0.130) (-0.017,-0.014) (-0.011,-0.040) (-0.005,-0.041) (-0.055,-0.038) (-0.010,-0.028)
 (-0.002,-0.019) (0.004,0.013) (0.005,0.002) (-0.001,0.023) (0.098,0.038) (-0.058,-0.009) (-0.008,-0.010) (-0.049,0.011) (0.035,-0.040) (-0.008,0.023) (0.010,0.037) (0.009,0.035)
 (-0.010,-0.009) (-0.011,-0.096) (0.003,0.016) (-0.002,-0.064) (0.004,0.015) (0.034,0.025) (-0.003,-0.003) (-0.004,0.029) (0.038,0.020) (0.002,0.006) (-0.008,-0.006) (-0.008,-0.047)
 (-0.004,-0.000) (-0.001,-0.014) (0.008,0.033) (-0.025,0.013) (0.006,0.015) (-0.046,0.004) (-0.001,0.001) (-0.007,-0.010) (-0.007,-0.016) (0.062,0.054) (-0.003,-0.016) (-0.014,0.003)
 (0.018,0.028) (0.000,-0.014) (-0.011,-0.005) (-0.003,0.031) (-0.010,0.016) (0.211,0.079) (0.011,-0.078) (0.001,-0.049) (0.114,0.139) (0.048,-0.002) (0.019,0.006) (-0.017,-0.001)
 (-0.000,-0.003) (-0.044,-0.014) (-0.005,0.007) (-0.025,-0.020) (0.031,0.030) (0.006,0.009) (0.010,0.054) (0.014,-0.073) (-0.003,-0.007) (-0.000,-0.161) (0.002,-0.013) (-0.063,-0.022)
 (-0.000,-0.023) (-0.002,-0.025) (0.001,-0.003) (-0.006,-0.010) (-0.010,-0.028) (-0.009,-0.008) (-0.000,-0.012) (0.022,0.052) (-0.004,0.003) (-0.005,-0.047) (-0.000,0.029) (0.006,-0.019)
 (-0.005,0.003) (-0.007,0.012) (-0.009,0.044) (-0.002,0.012) (-0.008,-0.042) (0.053,0.010) (0.039,0.050) (0.024,-0.114) (0.035,-0.040) (0.047,0.045) (0.083,0.021) (-0.015,-0.013)
 (0.026,0.031) (-0.010,0.038) (0.008,-0.012) (0.002,-0.018) (0.021,0.028) (0.003,-0.019) (-0.033,0.038) (-0.014,0.007) (-0.018,-0.021) (-0.008,-0.014) (-0.022,-0.022) (0.019,0.009)
 (-0.012,0.000) (-0.019,-0.013) (0.000,-0.002) (-0.042,-0.002) (0.007,-0.024) (-0.018,-0.041) (-0.005,0.054) (-0.013,-0.051) (0.141,0.148) (-0.010,-0.056) (-0.106,-0.037) (0.052,0.034)
 (0.017,0.042) (-0.008,0.019) (-0.289,-0.144) (0.011,0.016) (0.061,0.016) (-0.008,0.022) (0.011,0.022) (0.003,0.008) (0.002,0.051) (0.000,0.002) (-0.004,-0.096) (0.006,0.016)
 (0.025,-0.032) (-0.003,0.005) (-0.022,0.018) (-0.017,0.008) (-0.019,0.014) (-0.070,-0.023) (-0.049,-0.105) (0.063,0.042) (0.003,-0.051) (-0.010,-0.005) (0.045,0.097) (-0.017,-0.003)
 (0.022,0.036) (0.011,-0.014) (-0.042,-0.005) (-0.116,-0.067) (-0.011,-0.015) (-0.011,0.020) (-0.003,-0.002) (0.161,0.071) (0.063,0.055) (0.062,0.058) (0.015,0.008) (-0.025,0.004)
 (-0.108,-0.068) (0.013,-0.042) (0.044,0.064) (-0.013,0.001) (-0.003,0.015) (0.032,0.012) (-0.010,0.006) (-0.020,0.012) (-0.016,0.035) (-0.008,0.018) (-0.010,-0.004) (0.037,0.059)
 (-0.004,0.003) (0.005,0.012) (0.047,0.069) (0.007,0.004) (0.006,0.006) (-0.010,0.006) (0.001,-0.012) (0.124,0.123) (0.025,-0.046) (0.001,-0.011) (0.145,0.166) (0.089,0.049)
 (-0.017,-0.010) (0.007,0.018) (-0.007,0.002) (-0.015,-0.002) (-0.011,-0.021) (-0.016,0.006) (-0.005,-0.005) (-0.004,-0.007) (0.010,-0.033) (0.054,0.108) (-0.004,-0.007) (-0.004,-0.094)
 (0.000,-0.023) (-0.023,-0.011) (0.006,-0.029) (-0.009,0.026) (-0.011,0.016) (-0.028,0.007) (-0.010,0.034) (-0.010,-0.105) (0.011,0.016) (0.020,-0.035) (0.085,0.067) (0.013,-0.010)
 (-0.005,-0.007) (0.002,0.004) (0.005,-0.003) (-0.008,-0.003) (-0.011,0.016) (-0.001,0.000) (-0.001,-0.019) (-0.000,0.016) (-0.031,-0.038) (-0.020,0.041) (-0.051,0.023) (-0.011,0.002)
 (-0.012,0.006) (-0.006,-0.001) (0.011,0.018) (0.080,0.050) (-0.005,0.011) (-0.011,0.020) (0.013,0.027) (-0.030,-0.005) (-0.062,0.000) (-0.018,-0.081) (-0.006,0.005) (-0.014,0.000)
 (-0.011,-0.013) (-0.009,-0.021) (-0.001,-0.005) (0.016,0.003) (-0.005,-0.015) (-0.034,-0.015) (0.007,-0.027) (-0.011,-0.025) (-0.046,0.001) (0.010,-0.017) (0.092,0.106) (0.021,-0.079)
 (-0.149,-0.054) (-0.031,-0.015) (0.002,-0.009) (-0.014,-0.015) (-0.008,-0.041) (-0.013,-0.042) (-0.018,0.015) (-0.018,-0.031) (-0.007,-0.013) (0.005,-0.081) (0.006,-0.039) (-0.005,-0.021)
 (-0.014,-0.165) (-0.002,-0.032) (-0.010,0.034) (-0.002,-0.024) (-0.008,0.004) (-0.007,-0.005) (-0.006,-0.007) (-0.031,0.013) (-0.003,0.004) (-0.032,-0.022) (-0.027,-0.051) (-0.005,-0.011)
 (0.014,0.034) (0.002,-0.022) (0.035,0.017) (-0.008,-0.011) (-0.016,-0.009) (-0.040,-0.028) (-0.008,-0.013) (0.008,-0.044) (-0.006,-0.010) (0.091,0.112) (0.003,0.046) (0.055,0.076)
 (0.001,0.010) (0.032,0.018) (-0.006,0.015) (-0.013,-0.006) (-0.016,0.000) (0.019,0.001) (-0.012,0.001) (-0.024,-0.041) (-0.008,-0.011) (-0.006,0.026) (-0.004,-0.026) (0.009,0.012)
 (0.022,-0.009) (0.013,0.019) (-0.003,-0.016) (-0.007,0.004) (-0.001,0.008) (-0.008,0.002) (-0.020,-0.000) (-0.007,0.015) (-0.005,0.005) (0.057,0.059) (0.007,-0.031) (-0.006,-0.060)
 (0.008,-0.069) (0.053,0.019) (-0.008,0.003) (0.033,0.003) (-0.009,-0.000) (-0.110,-0.086) (-0.009,-0.003) (0.025,-0.034) (-0.006,-0.007) (-0.001,-0.003) (0.001,-0.011) (-0.001,-0.013)
 (-0.013,0.012) (0.014,-0.019) (0.008,-0.041) (0.011,0.049) (0.005,-0.000) (-0.001,0.018) (0.002,0.001) (-0.003,0.010) (-0.001,-0.001) (-0.006,-0.016) (-0.005,-0.016) (0.008,0.040)
 (0.003,0.007) (0.020,0.029) (-0.006,-0.013) (0.021,0.036) (-0.019,-0.010) (-0.018,-0.000) (-0.020,-0.005) (-0.011,-0.004) (-0.006,-0.037) (-0.004,0.021) (0.022,0.080) (-0.022,0.092)
 (-0.026,0.093) (0.022,0.028) (0.014,0.047) (-0.006,-0.037) (0.012,0.026) (0.049,0.058) (0.001,0.014) (-0.008,-0.009) (0.003,0.018) (-0.014,0.009) (-0.059,-0.010) (-0.017,-0.045)
 (0.003,-0.003) (-0.002,-0.005) (0.009,0.039) (0.002,-0.012) (0.012,-0.001) (-0.000,0.046) (-0.002,-0.011) (-0.006,-0.005) (0.004,0.001) (-0.010,0.007) (-0.003,-0.011) (0.002,0.004)
 (0.009,0.007) (-0.007,0.036) (-0.001,0.007) (0.054,0.015) (-0.013,0.005) (0.018,-0.004) (0.036,0.054) (-0.006,0.015) (0.016,0.012) (0.023,0.003) (-0.003,0.003) (-0.012,0.159)
 (0.206,0.314) (-0.011,-0.036) (-0.013,-0.036) (-0.015,-0.030) (0.048,0.077) (-0.009,-0.015) (-0.006,0.020) (-0.007,-0.009) (-0.009,-0.014) (-0.023,-0.019) (-0.023,-0.105) (-0.005,-0.005)
 (-0.009,0.004) (-0.004,0.001) (0.003,-0.006) (0.002,-0.001) (0.019,0.005) (0.001,0.001) (-0.011,-0.005) (-0.040,0.001) (-0.013,-0.004) (-0.001,-0.051) (-0.007,-0.005) (0.157,0.058)
 (0.039,0.028) (0.060,0.038) (0.004,0.001) (0.008,-0.011) (-0.068,-0.046) (-0.003,-0.028) (-0.001,0.009) (0.013,0.007) (-0.005,-0.006) (-0.011,-0.014) (-0.002,0.011) (-0.006,-0.009)
 (0.005,0.059) (-0.020,-0.012) (-0.008,0.012) (-0.016,-0.018) (-0.010,-0.005) (-0.011,0.004) (0.040,-0.011) (-0.002,0.001) (-0.017,-0.017) (-0.006,-0.020) (-0.005,-0.008) (-0.034,0.015)
 (0.007,-0.006) (0.038,0.074) (-0.031,0.002) (-0.061,-0.016) (-0.029,0.001) (-0.008,0.000) (-0.013,0.005) (-0.008,-0.029) (-0.013,-0.016) (0.016,0.011) (-0.020,-0.000) (-0.013,0.004)
 (0.004,0.023) (0.072,0.184) (-0.003,-0.003) (-0.009,-0.028) (-0.015,-0.011) (-0.009,-0.028) (-0.006,-0.029) (-0.005,-0.008) (-0.008,-0.010) (-0.002,-0.010) (-0.012,-0.014) (-0.005,-0.029)
 (-0.019,0.014) (-0.004,0.031) (0.129,0.050) (0.068,0.063) (-0.010,0.027) (-0.064,0.029) (0.016,0.026) (0.000,0.008) (0.061,-0.020) (-0.000,0.024) (-0.014,-0.166) (-0.002,-0.010)
 (0.099,0.061) (0.023,0.042) (0.038,0.053) (-0.001,-0.008) (0.009,-0.004) (-0.027,0.005) (-0.008,0.040) (-0.002,-0.009) (-0.016,-0.024) (-0.042,-0.041) (0.026,0.027) (-0.026,0.002)
 (0.051,0.022) (-0.012,-0.022) (-0.044,-0.027) (-0.024,-0.023) (-0.002,0.015) (0.012,0.020) (0.022,-0.002) (-0.006,-0.011) (0.007,0.022) (-0.008,-0.001) (-0.006,0.020) (0.001,-0.005)
 (0.103,0.087) (0.014,-0.037) (0.006,-0.019) (0.057,0.057) (0.061,0.022) (-0.022,-0.015) (-0.005,0.001) (-0.002,-0.010) (-0.068,-0.042) (0.028,0.008) (-0.027,-0.007) (-0.004,0.035)
 (0.004,0.003) (0.003,-0.049) (0.040,0.048) (-0.005,-0.013) (-0.006,-0.005) (-0.009,0.004) (-0.008,-0.038) (-0.007,-0.013) (-0.011,0.023) (-0.007,-0.005) (-0.001,0.008) (-0.059,-0.035)
 (-0.014,-0.042) (0.052,0.035) (0.005,-0.021) (-0.028,-0.049) (0.042,0.086) (-0.017,0.020) (-0.001,0.042) (0.025,0.001) (-0.057,-0.011) (-0.113,-0.090) (-0.011,-0.009) (-0.017,0.074)
 (-0.014,0.008) (0.003,0.015) (-0.002,0.040) (0.004,0.058) (0.002,0.013) (-0.014,-0.015) (-0.004,-0.011) (0.004,-0.008) (-0.007,-0.010) (-0.003,-0.010) (0.035,0.044) (0.042,0.045)
 (0.015,0.027) (-0.016,-0.006) (-0.003,-0.006) (-0.009,-0.016) (-0.006,-0.022) (0.014,0.041) (0.007,-0.051) (-0.006,0.005) (-0.005,0.125) (0.004,-0.019) (-0.017,-0.004) (-0.009,-0.014)
 (-0.002,0.036) (-0.057,-0.024) (0.001,-0.019) (-0.079,-0.082) (-0.029,-0.015) (-0.013,0.017) (0.020,-0.002) (-0.049,-0.009) (-0.001,0.014) (-0.031,0.052) (-0.014,0.002) (0.004,-0.009)
 (-0.007,0.015) (-0.000,0.003) (-0.005,-0.059) (-0.003,-0.027) (-0.009,0.013) (0.102,0.067) (-0.010,-0.005) (0.082,0.014) (-0.016,-0.010) (-0.007,0.006) (-0.010,-0.006) (-0.009,-0.002)
 (-0.017,-0.001) (-0.003,-0.038) (0.005,0.007) (-0.005,-0.014) (0.042,0.044) (-0.010,-0.012) (-0.003,-0.002) (-0.007,0.009) (0.003,0.014) (0.086,0.004) (-0.048,-0.034) (-0.014,-0.007)
 (-0.006,-0.098) (-0.005,-0.010) (0.057,-0.001) (0.008,0.020) (0.019,0.012) (-0.007,-0.003) (0.010,-0.016) (-0.010,-0.007) (0.008,0.020) (0.057,0.050) (0.008,-0.020) (-0.003,-0.014)
 (-0.037,0.006) (0.007,-0.018) (-0.038,0.017) (-0.012,0.058) (-0.005,-0.015) (-0.005,-0.001) (-0.006,0.008) (-0.003,0.009) (-0.001,-0.032) (-0.003,-0.012) (-0.002,0.016) (-0.016,-0.010)
 (-0.002,0.004) (-0.012,-0.037) (0.007,0.015) (-0.008,-0.015) (-0.004,0.025) (-0.015,0.053) (0.032,-0.004) (0.008,0.003) (-0.017,0.007) (-0.009,-0.000) (-0.215,-0.102) (0.006,0.013)
 (0.069,0.020) (-0.008,0.029) (0.007,0.009) (0.001,0.019) (0.059,0.202) (-0.008,-0.002) (-0.011,0.027) (-0.004,-0.025) (0.013,0.052) (0.001,-0.019) (-0.006,0.010) (-0.007,0.006)
 (-0.006,0.003) (-0.019,0.018) (-0.005,-0.003) (0.008,-0.001) (0.022,-0.001) (-0.018,0.017) (0.009,0.038) (-0.001,-0.010) (0.043,0.029) (-0.014,-0.022) (-0.016,-0.002) (-0.046,-0.017)
 (-0.011,-0.015) (0.004,-0.001) (0.000,0.024) (0.019,-0.011) (0.001,0.037) (-0.000,0.025) (0.006,-0.006) (-0.007,0.002) (-0.012,-0.004) (0.006,-0.002) (0.001,0.007) (-0.048,-0.020)
 (0.001,0.013) (0.047,0.007) (0.004,0.008) (-0.010,0.012) (0.002,0.026) (-0.038,-0.026) (0.021,0.015) (-0.005,-0.006) (0.005,-0.006) (-0.004,0.010) (-0.000,-0.013) (0.002,0.002)
 (-0.003,0.004) (-0.009,-0.008) (0.006,0.026) (-0.049,-0.020) (-0.015,-0.027) (-0.055,-0.063) (-0.017,0.014) (-0.021,0.013) (-0.029,-0.014) (0.000,0.010) (-0.002,0.006) (0.145,0.033)
 (0.009,-0.010) (-0.057,-0.025) (0.005,0.004) (0.006,-0.007) (0.004,-0.105) (-0.000,-0.019) (0.002,-0.013) (0.017,0.042) (0.003,-0.035) (0.069,0.037) (0.001,-0.003) (-0.009,0.012)
 (-0.004,-0.000) (-0.006,-0.003) (-0.039,0.002) (-0.014,0.082) (0.009,0.009) (0.002,0.027) (0.015,0.004) (0.061,0.031) (-0.003,-0.031) (0.030,0.023) (-0.020,-0.016) (-0.021,-0.011)
 (-0.038,-0.018) (-0.004,-0.015) (0.007,-0.037) (0.019,0.026) (0.015,0.047) (0.012,0.051) (0.010,0.043) (0.024,0.034) (0.022,0.049) (-0.002,-0.023) (0.012,0.071) (0.033,0.066)
 (0.025,-0.020) (-0.003,-0.034) (-0.033,0.053) (0.004,-0.022) (-0.096,-0.010) (-0.020,-0.029) (-0.100,-0.068) (-0.002,0.010) (0.010,0.014) (0.003,0.017) (0.005,0.006) (-0.011,-0.120)
 (-0.007,0.014) (-0.000,0.013) (-0.012,0.013) (-0.015,0.004) (0.025,-0.010) (-0.006,-0.036) (0.029,0.077) (-0.032,-0.014) (0.063,-0.006) (0.067,-0.007) (0.001,-0.005) (-0.007,-0.005)
 (-0.074,-0.009) (-0.014,0.006) (0.090,0.054) (-0.014,-0.025) (-0.011,-0.013) (0.010,-0.089) (0.025,0.057) (-0.013,-0.005) (-0.006,-0.102) (0.005,-0.022) (-0.047,-0.018) (-0.005,-0.032)
 (-0.018,-0.025) (-0.001,0.005) (-0.012,-0.001) (-0.039,-0.013) (0.002,-0.077) (0.007,0.029) (0.002,-0.007) (0.066,0.042) (0.002,0.013) (-0.002,0.019) (0.001,-0.003) (0.004,0.023)
 (-0.008,-0.009) (-0.014,-0.026) (-0.001,0.007) (-0.004,-0.002) (-0.001,-0.002) (0.004,-0.014) (0.006,-0.106) (0.008,-0.069) (-0.000,-0.015) (-0.000,-0.037) (0.000,0.010) (-0.001,0.032)
 (0.009,-0.037) (0.000,-0.004) (-0.007,-0.030) (-0.003,-0.030) (-0.010,0.011) (0.021,0.060) (-0.003,0.296) (0.007,-0.010) (0.001,-0.014) (0.014,0.042) (0.006,-0.021) (0.008,0.009)
 (0.013,0.043) (0.005,-0.015) (-0.004,-0.009) (-0.020,0.022) (-0.001,0.029) (0.053,-0.030) (-0.006,-0.009) (-0.068,-0.024) (-0.057,0.054) (0.150,0.051) (0.000,-0.018) (-0.003,-0.001)
 (-0.001,-0.012) (-0.010,-0.037) (-0.005,-0.006) (0.011,0.011) (-0.001,-0.021) (0.003,-0.010) (0.006,0.051) (0.016,-0.000) (0.012,0.028) (-0.007,-0.074) (-0.005,-0.006) (-0.006,0.002)
 (-0.006,-0.000) (-0.020,-0.020) (-0.005,-0.011) (-0.013,-0.013) (-0.049,-0.026) (0.004,-0.134) (0.024,0.046) (-0.006,0.025) (-0.007,0.003) (0.001,-0.038) (0.007,-0.007) (0.000,-0.041)
 (-0.027,0.001) (0.068,0.040) (0.105,0.075) (0.004,0.000) (0.003,-0.111) (-0.000,0.002) (0.036,0.013) (-0.006,0.013) (0.003,0.011) (0.030,0.009) (-0.025,0.015) (-0.009,0.002)
 (0.002,-0.054) (0.044,0.040) (0.020,0.007) (-0.004,-0.009) (-0.036,0.051) (-0.012,-0.001) (-0.023,0.060) (-0.010,0.083) (-0.048,-0.008) (0.012,0.023) (-0.014,0.006) (-0.003,-0.028)
 (-0.005,0.006) (0.034,0.035) (-0.005,-0.014) (-0.031,0.002) (-0.009,0.055) (-0.010,0.076) (0.047,-0.055) (-0.034,-0.029) (-0.011,-0.012) (0.072,0.160) (0.087,-0.005)};
\draw[coloncol!60,line width=.3pt] (axis cs:0.484,0.198) -- (axis cs:0.290,-0.080);
\addplot[only marks,mark=*,mark size=1.6pt,color=coloncol] coordinates { (0.484,0.198) };
\node[font=\sffamily\tiny,anchor=west,color=coloncol,inner sep=1pt] at (axis cs:0.300,-0.080) {\textit{IL7R}--CD127};
\draw[coloncol!60,line width=.3pt] (axis cs:0.629,0.124) -- (axis cs:0.290,-0.165);
\addplot[only marks,mark=*,mark size=1.6pt,color=coloncol] coordinates { (0.629,0.124) };
\node[font=\sffamily\tiny,anchor=west,color=coloncol,inner sep=1pt] at (axis cs:0.300,-0.165) {\textit{FCGR3A}--CD16};
\draw[coloncol!60,line width=.3pt] (axis cs:0.514,0.101) -- (axis cs:0.290,-0.250);
\addplot[only marks,mark=*,mark size=1.6pt,color=coloncol] coordinates { (0.514,0.101) };
\node[font=\sffamily\tiny,anchor=west,color=coloncol,inner sep=1pt] at (axis cs:0.300,-0.250) {\textit{IGHD}--IgD};
\draw[coloncol!60,line width=.3pt] (axis cs:0.515,0.066) -- (axis cs:0.290,-0.335);
\addplot[only marks,mark=*,mark size=1.6pt,color=coloncol] coordinates { (0.515,0.066) };
\node[font=\sffamily\tiny,anchor=west,color=coloncol,inner sep=1pt] at (axis cs:0.300,-0.335) {\textit{KLRB1}--CD161};
\end{axis}
\begin{axis}[at={(b.east)},anchor=west,xshift=1.3cm,name=c,width=.34\textwidth,height=.34\textwidth,xmin=-0.6,xmax=0.8,ymin=-0.6,ymax=0.8,
 xlabel={Held-out correlation},ylabel={Predicted correlation},title={\textbf{c}\enspace Within types, paired},
 axis x line*=bottom,axis y line*=left,tick align=outside,xmajorgrids=true,grid style={gray!18},axis line style={gray!55},tick label style={font=\sffamily\scriptsize},label style={font=\sffamily\small},clip=true,title style={font=\sffamily\small,at={(0,1)},anchor=base west,yshift=5pt}]
\draw[black!35] (axis cs:-0.6,-0.6) -- (axis cs:0.8,0.8);
\addplot[only marks,mark=*,mark size=.35pt,color=donorcol!45] coordinates {
 (-0.019,-0.016) (-0.014,-0.013) (-0.022,-0.018) (-0.054,-0.044) (-0.025,-0.013) (0.007,0.033) (0.005,0.009) (-0.008,-0.007) (-0.032,-0.034) (0.001,0.001) (-0.031,-0.036) (-0.001,0.004)
 (-0.025,-0.020) (0.018,0.035) (0.001,0.002) (-0.006,-0.004) (-0.122,-0.069) (0.080,0.120) (-0.024,-0.001) (-0.040,-0.002) (-0.030,-0.015) (-0.155,-0.131) (-0.015,-0.018) (0.014,0.024)
 (-0.008,0.003) (0.007,0.010) (-0.034,-0.046) (-0.095,-0.098) (-0.010,-0.035) (0.071,0.098) (-0.008,-0.006) (-0.224,-0.184) (-0.019,0.010) (-0.008,-0.016) (-0.001,-0.017) (-0.001,0.016)
 (-0.098,-0.074) (0.002,-0.004) (-0.001,-0.018) (-0.014,-0.017) (-0.010,-0.014) (0.025,0.035) (-0.005,-0.054) (0.026,0.094) (-0.035,-0.021) (0.065,0.077) (-0.029,-0.025) (-0.005,-0.018)
 (0.008,-0.006) (-0.047,-0.058) (-0.023,-0.023) (0.014,0.040) (0.002,-0.024) (-0.005,-0.015) (0.016,0.009) (-0.019,-0.003) (-0.186,-0.132) (0.002,0.004) (0.014,-0.002) (0.052,0.068)
 (0.006,0.015) (-0.025,-0.030) (-0.017,-0.030) (-0.021,-0.028) (-0.065,-0.102) (-0.057,-0.034) (0.083,0.096) (0.038,0.031) (-0.009,-0.020) (0.029,0.049) (-0.001,-0.005) (-0.023,-0.023)
 (0.015,0.009) (0.009,0.016) (0.039,0.045) (0.017,0.014) (-0.006,-0.002) (-0.002,-0.003) (-0.000,0.005) (-0.004,-0.005) (0.002,0.000) (-0.004,-0.003) (-0.005,-0.002) (-0.001,-0.002)
 (0.004,-0.015) (-0.016,-0.018) (0.005,-0.003) (0.078,0.103) (0.044,0.008) (0.021,0.073) (-0.065,-0.080) (-0.003,-0.004) (0.006,-0.019) (-0.021,-0.031) (0.080,0.186) (0.001,0.002)
 (-0.002,-0.016) (0.023,0.010) (-0.005,-0.010) (-0.004,0.010) (0.004,-0.009) (-0.008,-0.006) (-0.008,0.008) (-0.006,-0.011) (0.048,0.031) (0.020,0.013) (-0.036,-0.037) (-0.024,-0.024)
 (0.078,0.081) (-0.017,-0.017) (-0.459,-0.360) (-0.036,0.044) (-0.010,-0.021) (-0.009,-0.009) (0.014,0.024) (0.005,0.007) (-0.005,0.007) (-0.009,-0.021) (-0.002,-0.001) (0.001,0.005)
 (0.001,0.010) (-0.003,0.008) (-0.002,0.003) (-0.009,0.002) (-0.003,0.002) (-0.023,0.004) (-0.002,-0.004) (0.007,0.006) (0.003,-0.002) (0.002,0.012) (-0.062,-0.126) (-0.002,0.009)
 (-0.060,-0.072) (0.012,0.015) (0.016,0.040) (0.059,0.086) (0.009,0.010) (0.005,-0.007) (0.008,0.002) (-0.026,-0.026) (0.001,0.003) (0.003,-0.002) (-0.025,-0.032) (0.001,0.032)
 (-0.005,-0.013) (0.006,-0.006) (-0.030,-0.045) (0.003,0.024) (-0.038,-0.030) (0.086,0.062) (-0.004,-0.013) (-0.008,-0.010) (-0.002,-0.010) (-0.049,-0.042) (0.011,0.007) (0.000,0.005)
 (0.002,-0.002) (-0.002,-0.005) (-0.014,-0.027) (0.017,-0.005) (-0.009,-0.008) (0.009,-0.003) (-0.005,-0.009) (-0.001,-0.009) (-0.001,0.001) (0.058,0.060) (0.043,0.025) (-0.093,-0.081)
 (0.062,0.040) (-0.069,-0.043) (-0.014,-0.006) (-0.030,-0.066) (-0.006,0.009) (0.171,0.108) (-0.016,-0.016) (-0.023,0.003) (0.001,0.006) (-0.003,-0.001) (0.004,0.000) (-0.041,-0.027)
 (0.004,0.004) (-0.066,-0.054) (0.008,0.007) (-0.000,-0.007) (-0.008,-0.012) (0.017,0.021) (-0.060,-0.035) (-0.022,-0.018) (0.017,0.030) (-0.006,0.009) (-0.003,-0.016) (-0.037,-0.035)
 (-0.013,-0.017) (-0.016,-0.022) (0.038,0.027) (-0.059,-0.013) (-0.234,-0.148) (-0.022,-0.022) (-0.020,0.000) (-0.100,-0.059) (0.071,0.095) (-0.048,-0.016) (-0.053,-0.014) (-0.013,-0.008)
 (-0.126,-0.094) (0.001,-0.011) (0.017,0.055) (0.239,0.203) (0.006,0.008) (0.004,-0.022) (-0.075,0.010) (-0.033,-0.071) (-0.091,-0.034) (-0.007,-0.017) (-0.018,-0.047) (-0.014,-0.030)
 (-0.030,-0.024) (-0.000,-0.021) (-0.023,-0.013) (0.034,0.030) (-0.002,-0.013) (-0.052,-0.038) (-0.013,-0.013) (0.026,0.080) (-0.019,-0.052) (-0.034,-0.080) (-0.056,-0.084) (0.004,-0.015)
 (-0.008,-0.016) (-0.018,-0.017) (-0.009,-0.002) (-0.001,-0.001) (-0.030,-0.038) (-0.004,-0.004) (-0.063,-0.009) (-0.002,-0.006) (-0.001,0.005) (0.003,0.025) (-0.006,-0.015) (0.099,0.101)
 (-0.009,-0.014) (0.002,0.024) (-0.083,-0.102) (-0.017,-0.002) (0.001,-0.015) (0.002,-0.008) (-0.000,0.011) (-0.037,-0.056) (-0.006,-0.015) (0.010,0.002) (0.018,0.022) (0.003,0.002)
 (0.100,0.145) (-0.017,-0.015) (0.050,0.087) (0.019,0.000) (-0.014,-0.045) (-0.062,-0.067) (-0.013,-0.022) (-0.012,-0.027) (-0.003,-0.012) (0.235,0.179) (0.092,0.116) (0.105,0.088)
 (0.017,0.022) (0.000,-0.013) (-0.110,-0.090) (0.005,0.006) (-0.006,-0.015) (-0.020,-0.000) (-0.019,-0.009) (0.049,0.044) (-0.007,-0.016) (-0.021,-0.008) (-0.005,-0.012) (-0.314,-0.250)
 (-0.004,0.005) (0.001,-0.015) (-0.005,-0.031) (0.002,0.018) (-0.092,-0.059) (0.002,-0.017) (-0.000,-0.040) (-0.002,-0.002) (-0.005,0.001) (-0.005,-0.005) (-0.002,-0.010) (0.025,0.007)
 (-0.001,0.003) (-0.031,-0.033) (-0.018,-0.009) (-0.023,-0.019) (-0.031,-0.024) (0.004,-0.047) (-0.007,-0.019) (-0.025,-0.011) (-0.035,-0.030) (-0.014,-0.020) (-0.007,-0.013) (0.038,0.047)
 (0.008,-0.032) (-0.077,-0.079) (-0.005,-0.012) (-0.010,-0.016) (-0.002,-0.013) (0.002,-0.003) (0.006,-0.002) (0.005,0.013) (-0.016,-0.012) (-0.013,-0.011) (-0.020,-0.007) (-0.011,0.006)
 (0.001,-0.008) (0.124,0.176) (-0.022,-0.026) (-0.036,-0.020) (0.009,-0.003) (-0.010,-0.027) (0.029,0.011) (0.005,-0.015) (-0.013,-0.014) (-0.006,-0.017) (0.352,0.247) (0.079,0.112)
 (0.150,0.134) (0.002,0.002) (0.044,0.051) (-0.096,-0.073) (0.022,0.062) (0.017,0.063) (0.007,-0.007) (-0.013,-0.042) (0.077,0.156) (-0.010,-0.019) (-0.038,0.003) (-0.011,-0.037)
 (0.035,0.054) (-0.002,-0.007) (-0.005,-0.008) (0.001,-0.004) (-0.001,-0.004) (-0.001,-0.010) (0.010,-0.003) (-0.004,0.001) (-0.009,-0.008) (-0.001,-0.016) (0.109,0.109) (0.014,0.026)
 (0.006,-0.007) (0.149,0.135) (0.051,0.054) (-0.008,-0.005) (-0.008,0.004) (-0.001,0.004) (-0.076,-0.064) (-0.007,-0.012) (-0.009,-0.029) (0.001,0.007) (-0.001,0.003) (0.022,-0.004)
 (0.048,0.080) (0.013,-0.043) (-0.003,-0.014) (0.006,-0.016) (-0.126,-0.066) (-0.011,-0.019) (-0.002,0.005) (-0.000,0.003) (-0.003,-0.008) (-0.005,0.012) (0.000,0.001) (-0.003,0.005)
 (-0.008,-0.024) (-0.004,-0.004) (0.054,0.085) (-0.007,-0.008) (-0.026,-0.004) (-0.001,-0.002) (-0.005,-0.018) (0.028,0.006) (-0.004,-0.012) (-0.013,-0.013) (0.048,0.070) (0.001,0.039)
 (0.023,0.006) (0.040,0.027) (0.046,0.086) (0.078,0.123) (-0.031,-0.058) (0.001,-0.008) (-0.030,-0.035) (0.016,0.006) (-0.009,-0.014) (0.056,0.039) (-0.009,-0.013) (-0.001,0.001)
 (0.005,0.001) (-0.011,-0.005) (0.000,0.002) (0.002,-0.004) (0.001,0.009) (0.001,0.000) (-0.033,-0.018) (-0.000,0.006) (0.006,0.009) (-0.002,-0.022) (-0.010,-0.010) (0.071,0.067)
 (-0.044,-0.079) (0.023,0.085) (-0.047,-0.035) (0.143,0.134) (-0.025,-0.018) (0.045,0.059) (-0.019,-0.013) (-0.421,-0.333) (-0.025,0.025) (0.032,0.024) (-0.010,-0.011) (0.004,0.012)
 (-0.004,-0.020) (0.003,-0.004) (-0.005,-0.019) (0.029,0.035) (-0.003,-0.012) (0.058,0.046) (-0.003,-0.008) (-0.011,-0.002) (-0.015,-0.002) (-0.019,-0.009) (-0.057,0.071) (-0.018,-0.010)
 (-0.020,-0.035) (0.009,0.029) (-0.014,-0.014) (-0.027,-0.069) (0.004,0.006) (0.007,0.016) (-0.007,-0.010) (-0.023,-0.023) (-0.070,-0.049) (-0.018,-0.018) (0.002,0.002) (0.111,0.079)
 (0.088,0.043) (0.100,0.065) (0.102,0.061) (0.045,0.055) (0.126,0.111) (-0.083,-0.077) (-0.002,-0.021) (-0.011,-0.012) (-0.000,-0.003) (0.075,0.081) (-0.010,-0.024) (0.022,0.059)
 (-0.033,-0.049) (-0.007,0.005) (-0.060,-0.050) (-0.072,-0.038) (-0.005,-0.003) (-0.004,-0.000) (-0.004,0.000) (-0.045,-0.033) (-0.012,0.004) (0.000,-0.005) (-0.025,-0.030) (-0.008,0.031)
 (0.028,0.034) (-0.015,-0.075) (0.022,0.028) (-0.084,-0.033) (0.095,0.076) (0.071,0.044) (-0.078,-0.041) (-0.002,-0.001) (-0.034,-0.032) (0.090,0.030) (0.091,0.025) (-0.003,-0.007)
 (0.011,-0.000) (-0.009,-0.023) (0.044,0.045) (0.004,0.002) (-0.014,-0.022) (-0.012,-0.024) (0.099,0.082) (-0.007,-0.025) (-0.000,0.004) (0.001,0.005) (0.005,0.007) (0.010,0.026)
 (0.003,0.022) (-0.042,-0.047) (-0.015,-0.026) (0.514,0.361) (0.006,0.003) (-0.000,-0.008) (-0.001,-0.009) (0.021,0.006) (-0.017,0.002) (-0.034,-0.013) (-0.010,-0.011) (-0.009,-0.003)
 (-0.059,-0.016) (0.042,0.078) (-0.043,-0.015) (-0.037,0.004) (0.027,0.030) (-0.067,-0.020) (-0.005,-0.013) (-0.002,0.003) (-0.001,-0.003) (0.047,0.055) (-0.027,-0.029) (-0.002,-0.007)
 (-0.001,-0.011) (-0.032,-0.037) (-0.005,-0.012) (-0.000,-0.006) (-0.008,-0.007) (0.003,-0.022) (0.016,0.015) (-0.005,-0.007) (-0.072,-0.044) (0.001,-0.004) (-0.022,-0.005) (-0.006,-0.015)
 (-0.008,-0.016) (-0.049,-0.001) (-0.007,-0.033) (0.123,0.167) (-0.039,-0.020) (-0.070,-0.057) (-0.003,-0.004) (-0.003,-0.000) (-0.012,-0.004) (-0.025,-0.037) (-0.007,-0.003) (-0.016,0.024)
 (-0.003,-0.004) (-0.005,0.003) (-0.004,-0.002) (-0.003,0.007) (0.007,0.002) (0.012,0.015) (-0.001,0.008) (-0.007,-0.016) (0.003,0.017) (0.001,-0.002) (-0.005,-0.007) (0.011,0.011)
 (-0.054,-0.031) (-0.020,-0.019) (-0.001,0.018) (-0.007,0.009) (-0.005,-0.019) (-0.042,-0.066) (-0.011,0.013) (-0.079,-0.056) (0.014,0.009) (0.012,0.021) (0.104,0.132) (-0.000,0.009)
 (-0.014,-0.008) (-0.006,-0.018) (-0.024,-0.052) (-0.018,-0.022) (-0.040,-0.051) (-0.009,-0.018) (-0.016,-0.033) (-0.010,-0.010) (0.020,0.019) (0.026,0.023) (-0.021,-0.010) (-0.002,-0.011)
 (0.050,0.054) (-0.013,-0.002) (-0.000,-0.000) (-0.015,-0.016) (-0.005,-0.007) (-0.009,-0.010) (-0.007,-0.025) (0.016,0.004) (-0.012,-0.012) (-0.031,-0.027) (0.005,-0.007) (-0.042,-0.013)
 (-0.004,-0.007) (0.006,-0.008) (0.119,0.140) (-0.006,-0.095) (-0.000,-0.011) (0.047,0.047) (0.027,0.033) (-0.008,0.003) (0.007,0.003) (0.006,-0.003) (-0.005,-0.012) (0.112,0.048)
 (0.031,-0.003) (0.007,0.003) (0.063,0.051) (0.001,0.002) (0.088,0.033) (-0.003,-0.022) (-0.009,-0.010) (-0.000,-0.005) (-0.068,-0.018) (-0.008,-0.010) (-0.003,-0.006) (-0.000,-0.008)
 (-0.003,-0.001) (-0.006,-0.035) (-0.004,-0.006) (0.006,-0.004) (0.011,0.017) (-0.003,-0.002) (0.002,0.008) (-0.001,0.008) (0.015,0.015) (-0.008,-0.011) (0.062,0.067) (-0.045,0.010)
 (-0.007,0.002) (0.004,0.002) (0.022,0.048) (0.093,0.115) (0.006,0.024) (0.036,0.063) (0.062,0.069) (0.034,0.091) (-0.020,-0.034) (0.067,0.005) (-0.017,-0.006) (0.010,0.011)
 (0.021,0.011) (-0.079,-0.051) (0.030,-0.001) (-0.000,0.004) (-0.015,0.002) (-0.058,-0.096) (-0.027,-0.038) (0.005,0.064) (-0.017,-0.018) (-0.011,-0.008) (0.009,-0.003) (-0.005,-0.015)
 (-0.030,-0.041) (-0.017,-0.033) (-0.002,0.030) (0.027,0.036) (-0.020,-0.101) (0.074,0.096) (-0.045,-0.026) (0.019,0.017) (-0.016,-0.000) (0.003,0.006) (0.006,0.008) (0.050,0.053)
 (0.008,0.007) (0.002,-0.013) (-0.032,-0.021) (0.002,-0.000) (-0.022,-0.019) (0.060,0.015) (0.042,0.092) (0.029,0.013) (-0.024,-0.018) (0.202,0.152) (-0.011,-0.025) (-0.005,0.003)
 (-0.003,0.000) (-0.001,-0.003) (-0.001,-0.000) (0.004,0.004) (0.008,0.002) (-0.001,0.000) (0.028,0.009) (0.070,0.072) (-0.005,-0.016) (0.006,0.007) (-0.006,-0.010) (0.000,-0.010)
 (-0.009,-0.035) (-0.002,-0.015) (-0.004,-0.015) (-0.010,-0.020) (-0.047,-0.074) (-0.031,-0.034) (-0.050,-0.053) (-0.008,-0.029) (-0.006,-0.019) (-0.005,-0.010) (0.017,0.013) (-0.124,-0.084)
 (-0.027,-0.016) (-0.003,-0.007) (0.075,0.088) (0.008,0.035) (0.074,0.036) (-0.013,-0.008) (-0.009,-0.025) (-0.015,-0.019) (-0.011,-0.024) (0.017,0.006) (0.004,-0.016) (-0.050,-0.047)
 (-0.001,-0.005) (-0.041,-0.021) (-0.009,-0.009) (0.021,0.049) (-0.007,-0.025) (-0.010,-0.047) (-0.011,-0.026) (-0.009,-0.000) (-0.000,-0.005) (-0.010,0.002) (-0.001,0.007) (-0.003,0.005)
 (0.021,0.020) (-0.001,0.004) (-0.006,-0.015) (0.035,0.037) (-0.011,0.008) (0.006,-0.016) (0.122,0.079) (-0.017,-0.031) (-0.011,-0.013) (-0.005,-0.018) (-0.055,-0.003) (-0.010,-0.019)
 (-0.002,0.015) (0.004,0.009) (0.005,-0.005) (-0.001,0.045) (0.098,0.079) (-0.058,-0.035) (-0.008,-0.000) (-0.049,-0.047) (0.035,0.046) (-0.008,-0.005) (0.010,0.020) (0.009,0.008)
 (-0.010,-0.008) (-0.011,-0.010) (0.003,-0.003) (-0.002,-0.006) (0.004,0.028) (0.034,0.056) (-0.003,-0.002) (-0.004,0.016) (0.038,0.051) (0.002,0.035) (-0.008,-0.019) (-0.008,-0.016)
 (-0.004,-0.004) (-0.001,-0.001) (0.008,0.015) (-0.025,-0.037) (0.006,0.019) (-0.046,-0.034) (-0.001,0.001) (-0.007,-0.021) (-0.007,-0.007) (0.062,0.100) (-0.003,0.005) (-0.014,-0.004)
 (0.018,0.005) (0.000,-0.004) (-0.011,0.006) (-0.003,0.003) (-0.010,-0.013) (0.211,0.220) (0.011,-0.041) (0.001,-0.034) (0.114,0.114) (0.048,0.053) (0.019,0.013) (-0.017,-0.008)
 (-0.000,0.008) (-0.044,-0.036) (-0.005,-0.014) (-0.025,-0.053) (0.031,0.028) (0.006,0.002) (0.010,-0.002) (0.014,0.026) (-0.003,-0.017) (-0.000,-0.006) (0.002,-0.002) (-0.063,-0.036)
 (-0.000,-0.015) (-0.002,-0.006) (0.001,-0.001) (-0.006,-0.004) (-0.010,-0.010) (-0.009,-0.004) (-0.000,-0.004) (0.022,0.025) (-0.004,0.001) (-0.005,-0.018) (-0.000,-0.001) (0.006,0.009)
 (-0.005,-0.004) (-0.007,-0.003) (-0.009,0.003) (-0.002,0.004) (-0.008,0.005) (0.053,0.064) (0.039,0.047) (0.024,0.010) (0.035,0.035) (0.047,0.048) (0.083,0.104) (-0.015,-0.034)
 (0.026,0.025) (-0.010,0.022) (0.008,-0.015) (0.002,-0.014) (0.021,0.054) (0.003,0.006) (-0.033,-0.003) (-0.014,-0.012) (-0.018,-0.033) (-0.008,-0.021) (-0.022,-0.023) (0.019,-0.009)
 (-0.012,-0.014) (-0.019,-0.009) (0.000,-0.012) (-0.042,-0.031) (0.007,-0.006) (-0.018,-0.021) (-0.005,0.026) (-0.013,-0.060) (0.141,0.198) (-0.010,-0.026) (-0.106,-0.085) (0.052,0.018)
 (0.017,0.028) (-0.008,-0.009) (-0.289,-0.209) (0.011,0.025) (0.061,0.044) (-0.008,-0.016) (0.011,0.016) (0.003,-0.001) (0.002,-0.003) (0.000,-0.003) (-0.004,-0.005) (0.006,-0.002)
 (0.025,0.013) (-0.003,-0.001) (-0.022,-0.020) (-0.017,-0.020) (-0.019,-0.016) (-0.070,-0.063) (-0.049,-0.035) (0.063,0.076) (0.003,0.002) (-0.010,-0.008) (0.045,0.105) (-0.017,-0.022)
 (0.022,0.039) (0.011,0.008) (-0.042,-0.041) (-0.116,-0.130) (-0.011,-0.021) (-0.011,-0.026) (-0.003,-0.013) (0.161,0.124) (0.063,0.094) (0.062,0.060) (0.015,0.024) (-0.025,-0.032)
 (-0.108,-0.081) (0.013,0.009) (0.044,0.035) (-0.013,-0.006) (-0.003,-0.007) (0.032,0.037) (-0.010,-0.006) (-0.020,-0.007) (-0.016,-0.012) (-0.008,-0.007) (-0.010,0.005) (0.037,0.041)
 (-0.004,0.005) (0.005,0.005) (0.047,0.025) (0.007,0.002) (0.006,0.020) (-0.010,-0.010) (0.001,-0.014) (0.124,0.143) (0.025,-0.016) (0.001,-0.011) (0.145,0.145) (0.089,0.084)
 (-0.017,-0.011) (0.007,0.011) (-0.007,-0.009) (-0.015,-0.025) (-0.011,-0.006) (-0.016,-0.012) (-0.005,-0.004) (-0.004,-0.003) (0.010,-0.005) (0.054,0.042) (-0.004,-0.006) (-0.004,-0.004)
 (0.000,-0.005) (-0.023,-0.013) (0.006,-0.006) (-0.009,0.000) (-0.011,0.001) (-0.028,-0.008) (-0.010,0.014) (-0.010,-0.004) (0.011,0.001) (0.020,-0.010) (0.085,0.119) (0.013,0.003)
 (-0.005,0.001) (0.002,0.010) (0.005,-0.005) (-0.008,-0.007) (-0.011,0.007) (-0.001,-0.001) (-0.001,-0.002) (-0.000,-0.012) (-0.031,-0.058) (-0.020,-0.021) (-0.051,-0.051) (-0.011,-0.019)
 (-0.012,-0.034) (-0.006,-0.007) (0.011,0.024) (0.080,0.048) (-0.005,0.024) (-0.011,-0.011) (0.013,0.017) (-0.030,-0.014) (-0.062,-0.014) (-0.018,-0.029) (-0.006,-0.014) (-0.014,-0.024)
 (-0.011,-0.019) (-0.009,-0.005) (-0.001,-0.013) (0.016,0.004) (-0.005,0.000) (-0.034,-0.024) (0.007,0.001) (-0.011,-0.017) (-0.046,-0.041) (0.010,-0.004) (0.092,0.139) (0.021,0.005)
 (-0.149,-0.126) (-0.031,-0.011) (0.002,-0.006) (-0.014,-0.017) (-0.008,-0.060) (-0.013,-0.024) (-0.018,0.036) (-0.018,-0.016) (-0.007,-0.009) (0.005,-0.015) (0.006,0.044) (-0.005,-0.028)
 (-0.014,-0.012) (-0.002,-0.015) (-0.010,0.017) (-0.002,-0.025) (-0.008,0.018) (-0.007,0.005) (-0.006,-0.012) (-0.031,0.134) (-0.003,0.011) (-0.032,-0.076) (-0.027,-0.053) (-0.005,-0.029)
 (0.014,0.016) (0.002,-0.005) (0.035,0.036) (-0.008,-0.011) (-0.016,-0.015) (-0.040,-0.052) (-0.008,-0.000) (0.008,-0.002) (-0.006,-0.005) (0.091,0.057) (0.003,0.011) (0.055,0.038)
 (0.001,0.014) (0.032,0.061) (-0.006,-0.006) (-0.013,-0.004) (-0.016,-0.011) (0.019,0.011) (-0.012,-0.008) (-0.024,-0.031) (-0.008,-0.010) (-0.006,-0.005) (-0.004,-0.002) (0.009,0.006)
 (0.022,0.007) (0.013,0.062) (-0.003,-0.016) (-0.007,-0.016) (-0.001,-0.004) (-0.008,-0.008) (-0.020,0.005) (-0.007,-0.004) (-0.005,-0.010) (0.057,0.070) (0.007,-0.049) (-0.006,-0.028)
 (0.008,0.023) (0.053,0.037) (-0.008,-0.006) (0.033,0.018) (-0.009,-0.009) (-0.110,-0.097) (-0.009,0.019) (0.025,0.019) (-0.006,-0.004) (-0.001,0.016) (0.001,-0.009) (-0.001,-0.005)
 (-0.013,-0.009) (0.014,0.031) (0.008,0.001) (0.011,0.063) (0.005,-0.003) (-0.001,-0.006) (0.002,-0.006) (-0.003,-0.001) (-0.001,-0.021) (-0.006,-0.005) (-0.005,-0.004) (0.008,0.016)
 (0.003,-0.001) (0.020,0.032) (-0.006,-0.015) (0.021,0.032) (-0.019,-0.020) (-0.018,-0.018) (-0.020,-0.023) (-0.011,-0.017) (-0.006,-0.010) (-0.004,-0.002) (0.022,0.021) (-0.022,-0.024)
 (-0.026,-0.016) (0.022,0.010) (0.014,0.052) (-0.006,-0.018) (0.012,0.046) (0.049,0.054) (0.001,0.009) (-0.008,-0.024) (0.003,0.052) (-0.014,-0.004) (-0.059,-0.003) (-0.017,-0.027)
 (0.003,0.001) (-0.002,-0.001) (0.009,0.051) (0.002,-0.006) (0.012,0.002) (-0.000,0.001) (-0.002,-0.002) (-0.006,-0.003) (0.004,-0.000) (-0.010,-0.014) (-0.003,-0.004) (0.002,-0.024)
 (0.009,-0.006) (-0.007,-0.004) (-0.001,-0.001) (0.054,0.032) (-0.013,-0.013) (0.018,0.004) (0.036,0.029) (-0.006,0.006) (0.016,-0.016) (0.023,0.006) (-0.003,0.001) (-0.012,-0.009)
 (0.206,0.050) (-0.011,-0.013) (-0.013,-0.007) (-0.015,-0.010) (0.048,0.038) (-0.009,-0.003) (-0.006,-0.004) (-0.007,-0.004) (-0.009,-0.013) (-0.023,-0.024) (-0.023,-0.017) (-0.005,-0.004)
 (-0.009,0.009) (-0.004,-0.004) (0.003,-0.019) (0.002,0.004) (0.019,0.016) (0.001,-0.005) (-0.011,-0.001) (-0.040,-0.022) (-0.013,-0.004) (-0.001,-0.003) (-0.007,-0.005) (0.157,0.110)
 (0.039,0.059) (0.060,0.049) (0.004,0.009) (0.008,-0.001) (-0.068,-0.050) (-0.003,0.001) (-0.001,-0.003) (0.013,0.004) (-0.005,-0.010) (-0.011,-0.003) (-0.002,-0.004) (-0.006,-0.014)
 (0.005,0.002) (-0.020,-0.026) (-0.008,-0.010) (-0.016,-0.001) (-0.010,-0.016) (-0.011,-0.001) (0.040,0.034) (-0.002,-0.009) (-0.017,-0.015) (-0.006,-0.011) (-0.005,-0.002) (-0.034,-0.010)
 (0.007,-0.004) (0.038,0.073) (-0.031,-0.019) (-0.061,-0.045) (-0.029,-0.020) (-0.008,-0.006) (-0.013,-0.013) (-0.008,-0.033) (-0.013,-0.017) (0.016,0.004) (-0.020,-0.013) (-0.013,-0.007)
 (0.004,-0.013) (0.072,0.060) (-0.003,-0.030) (-0.009,-0.013) (-0.015,-0.017) (-0.009,0.013) (-0.006,-0.009) (-0.005,-0.002) (-0.008,-0.001) (-0.002,0.002) (-0.012,0.014) (-0.005,0.008)
 (-0.019,-0.015) (-0.004,-0.009) (0.129,0.117) (0.068,0.107) (-0.010,-0.011) (-0.064,-0.039) (0.016,0.012) (0.000,0.002) (0.061,0.061) (-0.000,0.002) (-0.014,-0.013) (-0.002,-0.008)
 (0.099,0.071) (0.023,0.027) (0.038,0.030) (-0.001,0.002) (0.009,0.024) (-0.027,-0.018) (-0.008,0.008) (-0.002,0.003) (-0.016,-0.025) (-0.042,-0.033) (0.026,0.021) (-0.026,-0.019)
 (0.051,0.067) (-0.012,-0.005) (-0.044,-0.027) (-0.024,-0.017) (-0.002,-0.004) (0.012,0.011) (0.022,0.020) (-0.006,0.009) (0.007,0.004) (-0.008,0.007) (-0.006,-0.004) (0.001,-0.017)
 (0.103,0.117) (0.014,-0.033) (0.006,-0.009) (0.057,0.064) (0.061,0.059) (-0.022,-0.015) (-0.005,0.007) (-0.002,-0.005) (-0.068,-0.060) (0.028,0.009) (-0.027,-0.030) (-0.004,0.007)
 (0.004,0.006) (0.003,-0.012) (0.040,0.054) (-0.005,-0.024) (-0.006,-0.009) (-0.009,-0.011) (-0.008,0.013) (-0.007,-0.017) (-0.011,-0.012) (-0.007,-0.013) (-0.001,0.011) (-0.059,-0.041)
 (-0.014,-0.021) (0.052,0.038) (0.005,0.010) (-0.028,-0.034) (0.042,0.122) (-0.017,-0.033) (-0.001,0.013) (0.025,0.007) (-0.057,-0.048) (-0.113,-0.129) (-0.011,-0.028) (-0.017,-0.024)
 (-0.014,-0.004) (0.003,-0.000) (-0.002,-0.003) (0.004,-0.003) (0.002,-0.001) (-0.014,-0.020) (-0.004,0.002) (0.004,-0.005) (-0.007,-0.007) (-0.003,-0.010) (0.035,0.040) (0.042,0.027)
 (0.015,0.036) (-0.016,-0.025) (-0.003,0.002) (-0.009,-0.010) (-0.006,-0.006) (0.014,0.043) (0.007,0.001) (-0.006,-0.010) (-0.005,-0.014) (0.004,-0.001) (-0.017,0.009) (-0.009,-0.001)
 (-0.002,0.015) (-0.057,-0.043) (0.001,-0.038) (-0.079,-0.103) (-0.029,-0.017) (-0.013,-0.016) (0.020,0.006) (-0.049,-0.015) (-0.001,-0.001) (-0.031,-0.018) (-0.014,0.001) (0.004,-0.009)
 (-0.007,-0.007) (-0.000,0.007) (-0.005,-0.006) (-0.003,-0.008) (-0.009,0.001) (0.102,0.069) (-0.010,-0.006) (0.082,0.050) (-0.016,-0.013) (-0.007,0.002) (-0.010,-0.007) (-0.009,0.001)
 (-0.017,0.003) (-0.003,0.005) (0.005,-0.002) (-0.005,-0.008) (0.042,0.042) (-0.010,0.001) (-0.003,-0.005) (-0.007,-0.006) (0.003,-0.005) (0.086,0.084) (-0.048,0.005) (-0.014,-0.008)
 (-0.006,0.002) (-0.005,-0.001) (0.057,0.034) (0.008,0.012) (0.019,0.009) (-0.007,-0.004) (0.010,0.011) (-0.010,-0.012) (0.008,0.030) (0.057,0.067) (0.008,-0.003) (-0.003,-0.013)
 (-0.037,0.006) (0.007,-0.024) (-0.038,-0.013) (-0.012,-0.008) (-0.005,-0.012) (-0.005,-0.006) (-0.006,-0.009) (-0.003,-0.012) (-0.001,0.030) (-0.003,0.019) (-0.002,-0.010) (-0.016,-0.007)
 (-0.002,-0.011) (-0.012,-0.007) (0.007,0.018) (-0.008,-0.024) (-0.004,0.021) (-0.015,-0.012) (0.032,0.038) (0.008,0.006) (-0.017,0.001) (-0.009,-0.009) (-0.215,-0.176) (0.006,0.031)
 (0.069,0.057) (-0.008,-0.011) (0.007,0.026) (0.001,-0.006) (0.059,0.022) (-0.008,-0.005) (-0.011,-0.006) (-0.004,-0.004) (0.013,0.022) (0.001,-0.004) (-0.006,0.000) (-0.007,0.003)
 (-0.006,-0.004) (-0.019,0.043) (-0.005,-0.004) (0.008,-0.004) (0.022,0.024) (-0.018,-0.019) (0.009,0.039) (-0.001,-0.012) (0.043,0.050) (-0.014,-0.021) (-0.016,-0.026) (-0.046,-0.058)
 (-0.011,-0.015) (0.004,0.004) (0.000,0.001) (0.019,0.012) (0.001,0.004) (-0.000,0.000) (0.006,0.004) (-0.007,0.001) (-0.012,-0.007) (0.006,-0.008) (0.001,0.005) (-0.048,-0.036)
 (0.001,-0.004) (0.047,0.051) (0.004,0.005) (-0.010,-0.020) (0.002,0.014) (-0.038,-0.038) (0.021,0.012) (-0.005,-0.004) (0.005,-0.008) (-0.004,0.009) (-0.000,0.010) (0.002,-0.008)
 (-0.003,-0.007) (-0.009,-0.007) (0.006,0.038) (-0.049,-0.037) (-0.015,-0.049) (-0.055,-0.057) (-0.017,-0.022) (-0.021,-0.015) (-0.029,-0.012) (0.000,-0.000) (-0.002,0.006) (0.145,0.108)
 (0.009,-0.007) (-0.057,-0.038) (0.005,0.006) (0.006,0.001) (0.004,-0.013) (-0.000,-0.007) (0.002,0.001) (0.017,0.024) (0.003,-0.005) (0.069,0.066) (0.001,-0.008) (-0.009,-0.005)
 (-0.004,-0.006) (-0.006,-0.004) (-0.039,-0.010) (-0.014,-0.005) (0.009,0.012) (0.002,0.008) (0.015,0.015) (0.061,0.047) (-0.003,-0.012) (0.030,0.023) (-0.020,-0.015) (-0.021,-0.021)
 (-0.038,-0.053) (-0.004,-0.008) (0.007,0.001) (0.019,0.017) (0.015,0.035) (0.012,0.004) (0.010,0.021) (0.024,0.028) (0.022,0.055) (-0.002,-0.006) (0.012,0.023) (0.033,0.077)
 (0.025,0.006) (-0.003,-0.018) (-0.033,0.038) (0.004,-0.014) (-0.096,-0.022) (-0.020,-0.016) (-0.100,-0.071) (-0.002,0.015) (0.010,-0.001) (0.003,-0.015) (0.005,0.017) (-0.011,-0.010)
 (-0.007,-0.012) (-0.000,-0.008) (-0.012,-0.017) (-0.015,-0.022) (0.025,0.030) (-0.006,-0.036) (0.029,0.081) (-0.032,-0.015) (0.063,0.065) (0.067,0.012) (0.001,0.005) (-0.007,-0.005)
 (-0.074,-0.043) (-0.014,-0.015) (0.090,0.064) (-0.014,-0.009) (-0.011,-0.011) (0.010,-0.007) (0.025,0.046) (-0.013,-0.021) (-0.006,-0.014) (0.005,-0.013) (-0.047,-0.021) (-0.005,-0.012)
 (-0.018,-0.005) (-0.001,-0.003) (-0.012,-0.011) (-0.039,-0.019) (0.002,0.014) (0.007,0.024) (0.002,0.005) (0.066,0.077) (0.002,-0.012) (-0.002,0.006) (0.001,0.000) (0.004,0.009)
 (-0.008,-0.004) (-0.014,-0.006) (-0.001,0.004) (-0.004,0.000) (-0.001,0.000) (0.004,-0.003) (0.006,0.004) (0.008,0.002) (-0.000,-0.007) (-0.000,0.006) (0.000,-0.002) (-0.001,0.007)
 (0.009,0.004) (0.000,-0.004) (-0.007,-0.006) (-0.003,-0.004) (-0.010,-0.007) (0.021,0.012) (-0.003,-0.002) (0.007,-0.006) (0.001,0.004) (0.014,0.029) (0.006,0.002) (0.008,-0.002)
 (0.013,0.012) (0.005,-0.001) (-0.004,0.003) (-0.020,-0.018) (-0.001,0.004) (0.053,0.042) (-0.006,-0.017) (-0.068,-0.051) (-0.057,-0.025) (0.150,0.150) (0.000,-0.006) (-0.003,-0.003)
 (-0.001,-0.001) (-0.010,-0.017) (-0.005,-0.006) (0.011,0.018) (-0.001,-0.010) (0.003,-0.005) (0.006,0.003) (0.016,0.015) (0.012,0.038) (-0.007,-0.013) (-0.005,0.001) (-0.006,-0.004)
 (-0.006,-0.008) (-0.020,-0.018) (-0.005,-0.012) (-0.013,-0.019) (-0.049,-0.041) (0.004,0.010) (0.024,0.026) (-0.006,-0.004) (-0.007,-0.009) (0.001,0.019) (0.007,0.016) (0.000,-0.001)
 (-0.027,-0.013) (0.068,0.094) (0.105,0.199) (0.004,0.019) (0.003,0.012) (-0.000,0.009) (0.036,0.032) (-0.006,-0.003) (0.003,0.004) (0.030,0.042) (-0.025,-0.015) (-0.009,-0.008)
 (0.002,0.009) (0.044,0.042) (0.020,0.016) (-0.004,-0.009) (-0.036,-0.011) (-0.012,-0.016) (-0.023,-0.004) (-0.010,-0.006) (-0.048,-0.029) (0.012,0.016) (-0.014,0.007) (-0.003,-0.017)
 (-0.005,-0.002) (0.034,0.023) (-0.005,-0.007) (-0.031,-0.020) (-0.009,-0.014) (-0.010,-0.008) (0.047,0.064) (-0.034,-0.014) (-0.011,-0.008) (0.072,0.069) (0.087,0.086)};
\draw[coloncol!60,line width=.3pt] (axis cs:0.629,0.523) -- (axis cs:0.290,-0.080);
\addplot[only marks,mark=*,mark size=1.6pt,color=coloncol] coordinates { (0.629,0.523) };
\node[font=\sffamily\tiny,anchor=west,color=coloncol,inner sep=1pt] at (axis cs:0.300,-0.080) {\textit{FCGR3A}--CD16};
\draw[coloncol!60,line width=.3pt] (axis cs:0.484,0.387) -- (axis cs:0.290,-0.165);
\addplot[only marks,mark=*,mark size=1.6pt,color=coloncol] coordinates { (0.484,0.387) };
\node[font=\sffamily\tiny,anchor=west,color=coloncol,inner sep=1pt] at (axis cs:0.300,-0.165) {\textit{IL7R}--CD127};
\draw[coloncol!60,line width=.3pt] (axis cs:0.514,0.361) -- (axis cs:0.290,-0.250);
\addplot[only marks,mark=*,mark size=1.6pt,color=coloncol] coordinates { (0.514,0.361) };
\node[font=\sffamily\tiny,anchor=west,color=coloncol,inner sep=1pt] at (axis cs:0.300,-0.250) {\textit{IGHD}--IgD};
\draw[coloncol!60,line width=.3pt] (axis cs:0.515,0.296) -- (axis cs:0.290,-0.335);
\addplot[only marks,mark=*,mark size=1.6pt,color=coloncol] coordinates { (0.515,0.296) };
\node[font=\sffamily\tiny,anchor=west,color=coloncol,inner sep=1pt] at (axis cs:0.300,-0.335) {\textit{KLRB1}--CD161};
\end{axis}
\end{tikzpicture}
}
\caption{\textbf{Gene--protein dependence between and within cell types in the confirmatory cohort.}
Held-out cross-correlations against predictions, averaged over the \XsHaoDonors{} donors, for every sixteenth of the \XsGenes{}~$\times$~\XsProteins{} gene--protein pairs (grey) and for labelled pairs (coloured). \textbf{a}, All cells, pairing-free channel. \textbf{b}, Within level-1 cell types, pairing-free channel. \textbf{c}, Within level-1 cell types, paired closed form. Red, the four genes with the largest held-out within-type correlation with their own protein; blue, two lineage pairs. The diagonal marks perfect prediction. The labelled pairs and averages are descriptive and were chosen after scoring.}
\label{sfig:entries}
\end{figure}

